\documentclass[11pt,a4paper]{article}

% ==================== TEMEL PAKETLER (pdflatex) ====================
\usepackage[T1]{fontenc}
\usepackage[utf8]{inputenc}
\usepackage{lmodern}
\usepackage[a4paper,margin=2cm,top=2.4cm,bottom=2.4cm]{geometry}
\usepackage{xcolor}
\usepackage{colortbl}
\usepackage{array}
\usepackage{booktabs}
\usepackage{longtable}
\usepackage{tabularx}
\usepackage{listings}
\usepackage{fancyhdr}
\usepackage{amsmath}
\usepackage{amssymb}
\usepackage{hyperref}

% ==================== RENK PALETİ ====================
\definecolor{anaMavi}{HTML}{132C4A}
\definecolor{vurguMavi}{HTML}{2E6DA4}
\definecolor{acikMavi}{HTML}{EAF2F9}
\definecolor{turuncu}{HTML}{D9701E}
\definecolor{yesil}{HTML}{2E7D32}
\definecolor{acikYesil}{HTML}{E8F5E9}
\definecolor{kirmizi}{HTML}{C0392B}
\definecolor{mor}{HTML}{6C3483}
\definecolor{acikMor}{HTML}{F3E9F7}
\definecolor{acikGri}{HTML}{F4F5F7}
\definecolor{ortaGri}{HTML}{6B7280}
\definecolor{koyuGri}{HTML}{2D2D2D}
\definecolor{kodArka}{HTML}{FBFBFA}

% ==================== HYPERREF ====================
\hypersetup{
  colorlinks=true,
  linkcolor={[HTML]{2E6DA4}},
  urlcolor={[HTML]{2E6DA4}},
  pdftitle={C++ ile Dear ImGui Rehberi},
  pdfauthor={Arayuz Rehberi},
  bookmarksnumbered=true
}

% ==================== BAŞLIK STİLLERİ (titlesec olmadan) ====================
\setcounter{secnumdepth}{0}
\newcounter{secc}
\newcounter{subsecc}[secc]
\newcounter{subsubsecc}[subsecc]
\newcommand{\gsec}[1]{\par\phantomsection\stepcounter{secc}%
  \setcounter{subsecc}{0}%
  \vspace{18pt}{\color{anaMavi}\Large\bfseries\thesecc.\hspace{0.5em}#1}\par
  \vspace{2pt}{\color{turuncu}\rule{\linewidth}{1.6pt}}\par\vspace{8pt}%
  \addcontentsline{toc}{section}{\protect\numberline{\thesecc}#1}\markboth{#1}{#1}}
\newcommand{\gsub}[1]{\par\phantomsection\stepcounter{subsecc}%
  \setcounter{subsubsecc}{0}%
  \vspace{11pt}{\color{vurguMavi}\large\bfseries\thesecc.\thesubsecc\hspace{0.5em}#1}\par\vspace{4pt}%
  \addcontentsline{toc}{subsection}{\protect\numberline{\thesecc.\thesubsecc}#1}}
\newcommand{\gsubsub}[1]{\par\phantomsection\stepcounter{subsubsecc}%
  \vspace{7pt}{\color{koyuGri}\normalsize\bfseries #1}\par\vspace{2pt}}

% Kısım (Part) başlığı
\newcommand{\gpart}[2]{\clearpage
  \begin{center}\vspace*{1.5cm}
  {\color{turuncu}\rule{0.7\linewidth}{1.5pt}}\\[10pt]
  {\color{ortaGri}\Large KISIM #1}\\[8pt]
  {\color{anaMavi}\Huge\bfseries #2}\\[10pt]
  {\color{turuncu}\rule{0.7\linewidth}{1.5pt}}
  \end{center}\vspace{1cm}}

% ==================== BAŞLIK / ALTBİLGİ ====================
\setlength{\headheight}{14pt}
\pagestyle{fancy}
\fancyhf{}
\renewcommand{\headrulewidth}{0.4pt}
\renewcommand{\footrulewidth}{0.4pt}
\fancyhead[L]{\small\color{ortaGri}C++ ile Dear ImGui}
\fancyhead[R]{\small\color{ortaGri}\leftmark}
\fancyfoot[C]{\small\color{ortaGri}\thepage}

% ==================== KOD LİSTELEME (C++) ====================
\lstdefinestyle{cppstyle}{
  language=C++,
  basicstyle=\ttfamily\small\color{koyuGri},
  keywordstyle=\color{vurguMavi}\bfseries,
  commentstyle=\color{yesil}\itshape,
  stringstyle=\color{turuncu},
  numberstyle=\tiny\color{ortaGri},
  numbers=left,
  numbersep=8pt,
  backgroundcolor=\color{kodArka},
  frame=single,
  framesep=6pt,
  rulecolor=\color{ortaGri!40},
  breaklines=true,
  showstringspaces=false,
  tabsize=4,
  columns=fullflexible,
  keepspaces=true,
  extendedchars=true,
  literate={ı}{{\i}}1{İ}{{\.I}}1{ş}{{\c s}}1{Ş}{{\c S}}1{ğ}{{\u g}}1{Ğ}{{\u G}}1{ç}{{\c c}}1{Ç}{{\c C}}1{ö}{{\"o}}1{Ö}{{\"O}}1{ü}{{\"u}}1{Ü}{{\"U}}1{—}{{-{}-}}2{–}{{-}}1{’}{{'}}1{“}{{"}}1{”}{{"}}1,
  morekeywords={ImGui,ImVec2,ImVec4,ImGuiIO,ImGuiStyle,ImGuiCol,ImGuiWindowFlags,
    ImGuiCond,ImDrawList,ImU32,ImFont,ImFontConfig,ImFontAtlas,ImGuiContext,
    ImGuiTableFlags,ImGuiTreeNodeFlags,ImGuiInputTextFlags,ImGuiComboFlags,
    ImGuiSelectableFlags,ImGuiTabBarFlags,ImGuiDockNodeFlags,ImGuiConfigFlags,
    ImGuiViewport,ImGuiInputTextCallbackData,ImGuiKey,ImGuiMouseButton,
    ImGuiSliderFlags,ImGuiColorEditFlags,ImTextureID,ImGuiListClipper,
    GLFWwindow,GLuint,uint32_t,int32_t,uint8_t,size_t,nullptr,constexpr,noexcept,
    override,final,auto,std,string,vector,uint,int64_t}
}
\lstset{style=cppstyle}

% Kabuk/terminal stili
\lstdefinestyle{shstyle}{
  basicstyle=\ttfamily\small\color{koyuGri},
  backgroundcolor=\color{koyuGri!6},
  frame=single,framesep=6pt,rulecolor=\color{ortaGri!40},
  breaklines=true,showstringspaces=false,tabsize=4,columns=fullflexible,
  keepspaces=true,
  literate={ı}{{\i}}1{İ}{{\.I}}1{ş}{{\c s}}1{Ş}{{\c S}}1{ğ}{{\u g}}1{Ğ}{{\u G}}1{ç}{{\c c}}1{Ç}{{\c C}}1{ö}{{\"o}}1{Ö}{{\"O}}1{ü}{{\"u}}1{Ü}{{\"U}}1{—}{{-{}-}}2{–}{{-}}1,
}

% CMake stili
\lstdefinestyle{cmakestyle}{
  basicstyle=\ttfamily\small\color{koyuGri},
  keywordstyle=\color{mor}\bfseries,
  commentstyle=\color{yesil}\itshape,
  stringstyle=\color{turuncu},
  backgroundcolor=\color{acikMor!35},
  frame=single,framesep=6pt,rulecolor=\color{mor!40},
  breaklines=true,showstringspaces=false,tabsize=4,columns=fullflexible,
  keepspaces=true,
  morekeywords={cmake_minimum_required,project,add_executable,add_subdirectory,
    target_link_libraries,target_include_directories,find_package,set,
    add_library,target_compile_features,FetchContent_Declare,FetchContent_MakeAvailable,
    include,option,if,endif,PRIVATE,PUBLIC,VERSION,REQUIRED},
  literate={ı}{{\i}}1{İ}{{\.I}}1{ş}{{\c s}}1{Ş}{{\c S}}1{ğ}{{\u g}}1{Ğ}{{\u G}}1{ç}{{\c c}}1{Ç}{{\c C}}1{ö}{{\"o}}1{Ö}{{\"O}}1{ü}{{\"u}}1{Ü}{{\"U}}1,
}

% ==================== ÖZEL KUTULAR (tcolorbox olmadan) ====================
\newcommand{\callout}[4]{\par\smallskip\noindent\begingroup
  \setlength{\fboxsep}{8pt}\setlength{\fboxrule}{0.8pt}%
  \fcolorbox{#1}{#2}{\parbox{\dimexpr\linewidth-2\fboxsep-2\fboxrule\relax}{%
    {\bfseries\color{#1}#3}\par\smallskip #4}}\endgroup\par\smallskip}
\newcommand{\bilgi}[1]{\callout{vurguMavi}{acikMavi}{Bilgi}{#1}}
\newcommand{\ipucu}[1]{\callout{yesil}{acikYesil}{İpucu}{#1}}
\newcommand{\uyari}[1]{\callout{kirmizi}{kirmizi!7}{Dikkat}{#1}}
\newcommand{\ornek}[2][Örnek]{\callout{ortaGri!90!black}{acikGri}{#1}{#2}}
\newcommand{\notk}[2][Not]{\callout{mor}{acikMor}{#1}{#2}}

% Yardımcı komutlar
\newcommand{\code}[1]{\texttt{\color{koyuGri}#1}}
\renewcommand{\arraystretch}{1.35}

% ==================== BAŞLANGIÇ ====================
\begin{document}

% ---------- KAPAK ----------
\begin{titlepage}
\centering
\vspace*{1.4cm}
{\color{turuncu}\rule{\textwidth}{2pt}}\\[8pt]
{\color{anaMavi}\Huge\bfseries C++ ile Dear ImGui}\\[10pt]
{\color{vurguMavi}\LARGE Immediate-Mode GUI ile Araç ve Arayüz Geliştirme Rehberi}\\[8pt]
{\color{ortaGri}\large Mimariden Widget'lara, Stilden Docking'e, Bol Projeli Kapsamlı Başvuru}\\[6pt]
{\color{turuncu}\rule{\textwidth}{2pt}}\\[1.4cm]

\begin{center}
\fcolorbox{vurguMavi}{acikMavi}{\parbox{0.86\textwidth}{\centering\large
Bu belge önce arayüz programlamanın \textbf{iki paradigmasını}
(immediate-mode ve retained-mode), sonra Dear ImGui'nin \textbf{nasıl çalıştığını}
(kare döngüsü, ID sistemi, DrawList), ardından \textbf{tüm önemli widget'ları} ve
en sonda \textbf{sıfırdan yazılmış çok sayıda tam projeyi} adım adım anlatır.
Tüm kod örnekleri İngilizce isimlendirmeyle, gerçek üretim stiline uygun
yazılmıştır ve olduğu gibi derlenebilir.\vspace{4pt}}}
\end{center}

\vspace{0.7cm}
\begin{center}
\fcolorbox{ortaGri!50}{acikGri}{\parbox{0.86\textwidth}{\small
\textbf{Kapsam --- Kısım 1 (Temeller):} ImGui nedir \textbullet\ immediate vs retained
\textbullet\ mimari \& kare döngüsü \textbullet\ kurulum \& CMake \textbullet\ GLFW+OpenGL3 backend
\textbullet\ ilk pencere \textbullet\ ana döngü anatomisi \textbullet\ IO \& DeltaTime.\vspace{3pt}

\textbf{Kapsam --- Kısım 2 (Widget'lar):} pencereler \& bayraklar \textbullet\ buton/checkbox/radio
\textbullet\ slider \& drag \textbullet\ InputText \textbullet\ combo \& listbox \textbullet\ renk seçici
\textbullet\ ağaç \& seçilebilir \textbullet\ tablolar \textbullet\ menü \& tab bar \textbullet\ plot \textbullet\ ID sistemi \& layout.\vspace{3pt}

\textbf{Kapsam --- Kısım 3 (İleri):} stil \& tema \textbullet\ font \& Türkçe karakter
\textbullet\ docking \& viewport \textbullet\ DrawList ile özel çizim \textbullet\ InputText callback
\textbullet\ klavye/fare \textbullet\ ListClipper ile performans \textbullet\ yerleşik demo.\vspace{3pt}

\textbf{Kapsam --- Kısım 4 (Projeler):} yeniden kullanılabilir App iskeleti
\textbullet\ Yapılacaklar (To-Do) uygulaması \textbullet\ Log konsolu \textbullet\ Canlı sistem monitörü
\textbullet\ Renk paleti tasarımcısı \textbullet\ Hex/bellek görüntüleyici \textbullet\ Basit çizim tuvali
\textbullet\ Oyun hata ayıklama HUD'u.\vspace{3pt}}}
\end{center}

\vfill
{\color{ortaGri}\large Hazırlanma Tarihi: 19 Temmuz 2026}\\[4pt]
{\color{ortaGri} C++17 \textbullet\ Dear ImGui 1.90+ \textbullet\ GLFW 3.x \textbullet\ OpenGL 3.3 \textbullet\ CMake}
\vspace*{0.4cm}
\end{titlepage}

% ---------- İÇİNDEKİLER ----------
\tableofcontents
\thispagestyle{empty}
\clearpage

% #####################################################################
\gpart{1}{Dear ImGui'nin Temelleri}
% #####################################################################

% =====================================================================
\gsec{Dear ImGui Nedir?}
% =====================================================================

Dear ImGui (kısaca ImGui, ``Immediate Mode Graphical User Interface''),
C++ ile yazılmış, bağımlılığı olmayan (dependency-free) bir arayüz
kütüphanesidir. Omar Cornut tarafından geliştirilir. Amacı; oyun motorları,
grafik uygulamaları, hata ayıklama araçları ve içerik editörleri için
\textbf{hızlıca üretilebilen, minimum durum (state) tutan} arayüzler
yaratmaktır. Bugün Unreal Engine'den Blender eklentilerine, NVIDIA
araçlarından bağımsız oyun stüdyolarına kadar sektörde çok yaygın kullanılır.

ImGui bir \emph{işletim sistemi} arayüzü (Qt, GTK gibi) değildir. Kendine ait
bir pencereyi işletim sisteminden açmaz; bunun yerine sizin zaten sahip
olduğunuz bir OpenGL/DirectX/Vulkan bağlamına (context) \textbf{üçgenler
çizerek} arayüz üretir. Bu yüzden ona genelde bir oyun ya da render döngüsünün
\emph{içine} yerleştiririz.

\gsub{Neye Yarar, Neye Yaramaz?}

\begin{center}
\begin{tabularx}{\textwidth}{@{}>{\raggedright\arraybackslash}X >{\raggedright\arraybackslash}X@{}}
\toprule
\rowcolor{acikYesil}\color{yesil}\textbf{İdeal olduğu yerler} & \color{kirmizi}\cellcolor{kirmizi!8}\textbf{Uygun olmadığı yerler}\\
\midrule
Oyun içi hata ayıklama panelleri, editörler & Son kullanıcıya sunulan cilalı masaüstü uygulaması (bazen yine de kullanılır)\\
Grafik/render araçları, shader denemeleri & Mobil dokunmatik odaklı tüketici uygulamaları\\
Bilimsel görselleştirme, veri denetleyicileri & Karmaşık, tema-ağırlıklı web benzeri UI'lar\\
Hızlı prototipleme, ``bir slider koyayım'' anları & Erişilebilirlik (a11y) zorunlu kurumsal yazılım\\
\bottomrule
\end{tabularx}
\end{center}

\bilgi{ImGui'nin sloganı ``bloat-free''dir. Bir değeri ekrana bağlamak için
gözlemci (observer), sinyal-slot, XML layout dosyası veya kod üreteci
\emph{yoktur}. Bir \code{if} bloğu kadar basittir: \code{if (ImGui::Button("Kaydet")) save();}}

\gsub{Kısa Tarihçe ve Ekosistem}

ImGui 2014'te açık kaynak oldu ve MIT lisanslıdır. Etrafında zengin bir
ekosistem büyüdü:

\begin{center}
\begin{tabularx}{\textwidth}{@{}>{\bfseries\color{anaMavi}}l >{\raggedright\arraybackslash}X@{}}
\toprule
\rowcolor{acikMavi}\color{anaMavi}\textbf{Bileşen} & \color{anaMavi}\textbf{Açıklama}\\
\midrule
Çekirdek (\code{imgui.h}) & Tüm widget'lar, düzen, girdi işleme. Platformdan bağımsız.\\
\rowcolor{acikGri} Backend'ler & \code{imgui\_impl\_glfw}, \code{imgui\_impl\_opengl3}, \code{\_dx11}, \code{\_vulkan}, \code{\_sdl2} vb. Girdiyi ve çizimi platforma bağlar.\\
ImPlot & Yüksek performanslı bilimsel grafik/çizim eklentisi.\\
\rowcolor{acikGri} ImGuizmo & 3B sahnede taşıma/döndürme/ölçekleme gizmoları.\\
imgui-node-editor & Düğüm (node) tabanlı grafik editörleri.\\
\bottomrule
\end{tabularx}
\end{center}

% =====================================================================
\gsec{Immediate Mode ve Retained Mode}
% =====================================================================

ImGui'yi gerçekten anlamanın anahtarı, onun ``immediate mode'' (anlık kip)
felsefesini kavramaktır. Qt, GTK, HTML/DOM, Android View gibi klasik arayüz
sistemleri \textbf{retained mode} (kalıcı kip) çalışır. İkisi arasındaki fark,
kütüphaneyle kurduğunuz ilişkinin tamamını belirler.

\gsub{Retained Mode: Nesneler Kalıcıdır}

Retained mode'da bir düğme oluşturursunuz ve o düğme nesnesi \emph{bellekte
yaşamaya devam eder}. Ona bir metin verir, bir tıklama olayı bağlar (callback)
ve sonra unutursunuz. Arayüz ağacı (widget tree) bir veri yapısı olarak
saklanır; siz onu değiştirdikçe kütüphane ekranı günceller.

\begin{lstlisting}
// Retained mode (kavramsal, Qt benzeri sozde-kod)
QPushButton* button = new QPushButton("Kaydet");   // nesne kalici olusturulur
button->setEnabled(false);                          // durumu kutuphane tutar
connect(button, &QPushButton::clicked, this, &App::onSave); // olay baglanir
layout->addWidget(button);                          // agaca eklenir
// ... buradan sonra 'button' bellekte yasar, olaylari bekler
\end{lstlisting}

\gsub{Immediate Mode: Her Karede Yeniden Söylersiniz}

ImGui'de düğme diye bir nesne \emph{yoktur}. Her karede (frame),
``şu anda ekranda bir Kaydet düğmesi olsun'' dersiniz. Fonksiyon o karede
düğmeyi çizer ve o karede tıklanıp tıklanmadığını \code{bool} olarak
\emph{hemen} döndürür. Bir sonraki kare, aynı satırı tekrar çağırırsınız.

\begin{lstlisting}
// Immediate mode (gercek ImGui)
if (ImGui::Button("Kaydet")) {   // hem cizer hem "bu kare tiklandi mi?" doner
    save();                       // tiklandiysa hemen burada is yapariz
}
// Sonraki kare bu satir yine calisir. Dugme icin saklanan bir nesne YOK.
\end{lstlisting}

\notk[Zihniyet değişimi]{Immediate mode'da arayüz, uygulama durumunuzun (state)
her kare yeniden çizilen bir \emph{görüntüsüdür}. ``Düğmeyi devre dışı bırak''
gibi bir komut yollamaz, sadece ``bu karede düğme devre dışı görünsün''
dersiniz: \code{ImGui::BeginDisabled(isBusy);} \dots\ \code{ImGui::EndDisabled();}.
Senkronizasyon problemi ortadan kalkar: gösterdiğiniz şey her zaman
verinizin o anki hâlidir.}

\gsub{İki Yaklaşımın Karşılaştırması}

\begin{center}
\begin{tabularx}{\textwidth}{@{}>{\raggedright\arraybackslash}p{0.24\textwidth} >{\raggedright\arraybackslash}X >{\raggedright\arraybackslash}X@{}}
\toprule
\rowcolor{acikMavi}\color{anaMavi}\textbf{Konu} & \color{anaMavi}\textbf{Retained (Qt/HTML)} & \color{anaMavi}\textbf{Immediate (ImGui)}\\
\midrule
Durum sahibi & Kütüphane widget ağacını tutar & \textbf{Siz} tutarsınız, UI sadece yansıtır\\
\rowcolor{acikGri} Olay modeli & Callback / sinyal-slot & Fonksiyon dönüş değeri (\code{bool})\\
Bellek & Widget başına kalıcı nesne & Kare başına geçici, minimum\\
\rowcolor{acikGri} Senkronizasyon & UI ile veri ayrışabilir (bug kaynağı) & Ayrışamaz; her kare yeniden okunur\\
Performans profili & Az widget'ta ucuz & Her kare tüm UI çalışır (çoğu araç için sorun değil)\\
\rowcolor{acikGri} En iyi olduğu yer & Statik, karmaşık son-kullanıcı UI & Dinamik, sık değişen araç/editör UI\\
\bottomrule
\end{tabularx}
\end{center}

\uyari{Immediate mode ``her kare tüm arayüzü yeniden inşa ediyorum'' demektir.
Bu kulağa pahalı gelir ama modern donanımda binlerce widget bile mikrosaniyeler
sürer, çünkü ImGui yalnızca üçgen listesi üretir. Yine de \emph{ağır hesapları}
(dosya okuma, ağ) UI kodunun içine koymayın; onları ayrı yapıp sonucu
gösterin.}

% =====================================================================
\gsec{ImGui'nin Mimarisi ve Kare Döngüsü}
% =====================================================================

ImGui üç katmandan oluşur ve bu katmanları ayırt etmek kurulumu anlamayı
kolaylaştırır. Sizin uygulamanız bu üç katmanı her karede belli bir sırayla
sürer.

\gsub{Üç Katman}

\begin{center}
\begin{tabularx}{\textwidth}{@{}>{\bfseries\color{anaMavi}}l >{\raggedright\arraybackslash}X@{}}
\toprule
\rowcolor{acikMavi}\color{anaMavi}\textbf{Katman} & \color{anaMavi}\textbf{Sorumluluk}\\
\midrule
1. Çekirdek ImGui & Widget mantığı, düzen, girdi durumu. \code{ImGui::Button}, \code{ImGui::Begin} burada. Çıktısı: bir \textbf{çizim komut listesi} (DrawData).\\
\rowcolor{acikGri} 2. Platform backend & Klavye, fare, pencere boyutu, zamanı ImGui'ye \emph{besler}. Örn. \code{imgui\_impl\_glfw}.\\
3. Renderer backend & ImGui'nin ürettiği üçgenleri GPU'ya çizer. Örn. \code{imgui\_impl\_opengl3}.\\
\bottomrule
\end{tabularx}
\end{center}

ImGui hiçbir zaman doğrudan ekrana çizmez. Bir karenin sonunda,
\code{ImGui::Render()} çağrısı bir \code{ImDrawData} nesnesi hazırlar: bu,
``şu köşe tamponunu şu dokuyla, şu kırpma dikdörtgeninde çiz'' diyen komutların
listesidir. Renderer backend bu listeyi alıp OpenGL çağrılarına çevirir.

\gsub{Bir Karenin Anatomisi}

Her kare tam olarak şu sırayı izler. Bu sırayı bozmak (örneğin
\code{NewFrame} çağırmadan widget kullanmak) çökmeye yol açar.

\begin{center}
\begin{tabularx}{\textwidth}{@{}>{\bfseries}r >{\raggedright\arraybackslash}X@{}}
\toprule
\rowcolor{acikMavi}\color{anaMavi}\textbf{Adım} & \color{anaMavi}\textbf{Yapılan iş}\\
\midrule
1 & Platform olaylarını işle (\code{glfwPollEvents}).\\
\rowcolor{acikGri} 2 & Backend'lerin \code{NewFrame}'i: yeni girdiyi ImGui'ye aktar.\\
3 & \code{ImGui::NewFrame()}: yeni kareyi başlat.\\
\rowcolor{acikGri} 4 & \textbf{Sizin UI kodunuz}: \code{Begin/End}, tüm widget çağrıları.\\
5 & \code{ImGui::Render()}: DrawData'yı üret.\\
\rowcolor{acikGri} 6 & Ekranı temizle (\code{glClear}).\\
7 & Renderer backend: DrawData'yı çiz.\\
\rowcolor{acikGri} 8 & Tamponları takasla (\code{glfwSwapBuffers}).\\
\bottomrule
\end{tabularx}
\end{center}

\bilgi{4. adım, yani sizin UI kodunuz, tamamen \code{NewFrame} ile
\code{Render} arasında olmak zorundadır. Bu ikisinin arasına yazdığınız her
\code{ImGui::} çağrısı o kareye ait UI'yı tanımlar. Dışına yazarsanız ImGui
``hangi karedeyim?'' diyemez ve assert ile durur.}


% =====================================================================
\gsec{Kurulum ve Proje Yapısı}
% =====================================================================

ImGui'nin bir ``kur'' komutu yoktur; kaynak dosyaları doğrudan projenize
eklersiniz. Bu, kütüphanenin felsefesinin bir parçasıdır: derlenmiş bir
kütüphane (\code{.lib}/\code{.so}) yerine \emph{birkaç .cpp dosyasını} kendi
projenizle beraber derlersiniz. Bu bölümde GLFW + OpenGL3 backend'i ile
modern bir CMake projesi kuracağız.

\gsub{Gerekli Dosyalar}

Depodan (\url{https://github.com/ocornut/imgui}) şu dosyalara ihtiyacınız var:

\begin{center}
\begin{tabularx}{\textwidth}{@{}>{\ttfamily\color{koyuGri}}l >{\raggedright\arraybackslash}X@{}}
\toprule
\rowcolor{acikMavi}\normalfont\color{anaMavi}\textbf{Dosya} & \color{anaMavi}\textbf{Rol}\\
\midrule
imgui.cpp / imgui.h & Çekirdek kütüphane\\
\rowcolor{acikGri} imgui\_draw.cpp & Çizim komutları ve DrawList\\
imgui\_tables.cpp & Tablo widget'ları\\
\rowcolor{acikGri} imgui\_widgets.cpp & Standart widget'lar\\
imgui\_demo.cpp & Demo penceresi (öğrenmek için altın değerinde)\\
\rowcolor{acikGri} backends/imgui\_impl\_glfw.cpp & GLFW platform backend\\
backends/imgui\_impl\_opengl3.cpp & OpenGL3 renderer backend\\
\bottomrule
\end{tabularx}
\end{center}

\ipucu{\code{imgui\_demo.cpp} içindeki \code{ImGui::ShowDemoWindow()}, neredeyse
her widget'ın canlı bir örneğini gösterir. Öğrenirken demo penceresini açın,
beğendiğiniz widget'ı bulun, sonra \code{imgui\_demo.cpp} içinde o metni aratıp
kodunu kopyalayın. Bu, ImGui öğrenmenin en hızlı yoludur.}

\gsub{Önerilen Klasör Düzeni}

\begin{lstlisting}[style=shstyle]
imgui_proje/
  CMakeLists.txt
  src/
    main.cpp
  external/
    imgui/            <- ocornut/imgui deposu (git submodule ideal)
    glfw/             <- glfw/glfw deposu (veya sistemden find_package)
\end{lstlisting}

\gsub{CMakeLists.txt --- FetchContent ile Otomatik İndirme}

En pratik yöntem, CMake'in \code{FetchContent} modülüyle ImGui ve GLFW'yi
otomatik indirmektir. Böylece kimse elle dosya kopyalamaz.

\begin{lstlisting}[style=cmakestyle]
cmake_minimum_required(VERSION 3.16)
project(imgui_proje LANGUAGES CXX)

set(CMAKE_CXX_STANDARD 17)
set(CMAKE_CXX_STANDARD_REQUIRED ON)

include(FetchContent)

# --- GLFW ---
FetchContent_Declare(glfw
  GIT_REPOSITORY https://github.com/glfw/glfw.git
  GIT_TAG        3.4)
set(GLFW_BUILD_DOCS OFF CACHE BOOL "" FORCE)
set(GLFW_BUILD_TESTS OFF CACHE BOOL "" FORCE)
set(GLFW_BUILD_EXAMPLES OFF CACHE BOOL "" FORCE)
FetchContent_MakeAvailable(glfw)

# --- Dear ImGui ---
FetchContent_Declare(imgui
  GIT_REPOSITORY https://github.com/ocornut/imgui.git
  GIT_TAG        v1.91.5)
FetchContent_MakeAvailable(imgui)

# ImGui bir CMake hedefi sunmaz; kaynaklari elle derleriz.
add_library(imgui STATIC
  ${imgui_SOURCE_DIR}/imgui.cpp
  ${imgui_SOURCE_DIR}/imgui_draw.cpp
  ${imgui_SOURCE_DIR}/imgui_tables.cpp
  ${imgui_SOURCE_DIR}/imgui_widgets.cpp
  ${imgui_SOURCE_DIR}/imgui_demo.cpp
  ${imgui_SOURCE_DIR}/backends/imgui_impl_glfw.cpp
  ${imgui_SOURCE_DIR}/backends/imgui_impl_opengl3.cpp)

target_include_directories(imgui PUBLIC
  ${imgui_SOURCE_DIR}
  ${imgui_SOURCE_DIR}/backends)

target_link_libraries(imgui PUBLIC glfw)

# --- Uygulamamiz ---
add_executable(app src/main.cpp)
target_link_libraries(app PRIVATE imgui)

# OpenGL'i bagla (Linux/macOS/Windows farkli)
find_package(OpenGL REQUIRED)
target_link_libraries(app PRIVATE OpenGL::GL)
\end{lstlisting}

\begin{lstlisting}[style=shstyle]
# Derleme
cmake -B build -DCMAKE_BUILD_TYPE=Release
cmake --build build -j
./build/app
\end{lstlisting}

\bilgi{ImGui, OpenGL fonksiyon işaretçilerini yüklemek için varsayılan olarak
kendi hafif yükleyicisini (\code{imgui\_impl\_opengl3\_loader.h}) kullanır;
bu yüzden GLAD/GLEW eklemenize gerek yoktur. Kendi çiziminizi (sahne) OpenGL
ile yapacaksanız yine de bir yükleyici istersiniz, ama sadece ImGui için
gerekmez.}

% =====================================================================
\gsec{İlk Program: Boş Bir ImGui Penceresi}
% =====================================================================

Şimdi en küçük tam programı yazalım. Bu program bir GLFW penceresi açar,
ImGui'yi başlatır ve içinde tek bir metin gösteren bir ImGui penceresi çizer.
Bu iskelet, rehberdeki her örneğin temelidir.

\gsub{Tam Kaynak Kod}

\begin{lstlisting}
// main.cpp -- en kucuk calisan ImGui uygulamasi
#include "imgui.h"
#include "imgui_impl_glfw.h"
#include "imgui_impl_opengl3.h"
#include <GLFW/glfw3.h>
#include <cstdio>

static void glfw_error_callback(int error, const char* description) {
    std::fprintf(stderr, "GLFW Hata %d: %s\n", error, description);
}

int main() {
    // 1) GLFW'yi baslat
    glfwSetErrorCallback(glfw_error_callback);
    if (!glfwInit()) return 1;

    // OpenGL 3.3 Core baglami iste
    const char* glsl_version = "#version 330";
    glfwWindowHint(GLFW_CONTEXT_VERSION_MAJOR, 3);
    glfwWindowHint(GLFW_CONTEXT_VERSION_MINOR, 3);
    glfwWindowHint(GLFW_OPENGL_PROFILE, GLFW_OPENGL_CORE_PROFILE);

    GLFWwindow* window = glfwCreateWindow(1280, 720,
                                          "ImGui Ilk Pencere", nullptr, nullptr);
    if (!window) { glfwTerminate(); return 1; }
    glfwMakeContextCurrent(window);
    glfwSwapInterval(1);   // VSync acik

    // 2) ImGui baglamini olustur
    IMGUI_CHECKVERSION();
    ImGui::CreateContext();
    ImGuiIO& io = ImGui::GetIO();
    io.ConfigFlags |= ImGuiConfigFlags_NavEnableKeyboard;  // klavye ile gezinme

    ImGui::StyleColorsDark();   // koyu tema

    // 3) Backend'leri baslat
    ImGui_ImplGlfw_InitForOpenGL(window, true);
    ImGui_ImplOpenGL3_Init(glsl_version);

    // 4) Ana dongu
    while (!glfwWindowShouldClose(window)) {
        glfwPollEvents();

        // Yeni kareyi baslat
        ImGui_ImplOpenGL3_NewFrame();
        ImGui_ImplGlfw_NewFrame();
        ImGui::NewFrame();

        // ---- BURADAN itibaren bizim UI kodumuz ----
        ImGui::Begin("Merhaba");
        ImGui::Text("Merhaba, Dear ImGui!");
        ImGui::Text("Uygulama %.1f FPS calisiyor", io.Framerate);
        ImGui::End();
        // ---- UI kodu bitti ----

        // Kareyi render et
        ImGui::Render();
        int display_w, display_h;
        glfwGetFramebufferSize(window, &display_w, &display_h);
        glViewport(0, 0, display_w, display_h);
        glClearColor(0.10f, 0.11f, 0.13f, 1.0f);
        glClear(GL_COLOR_BUFFER_BIT);
        ImGui_ImplOpenGL3_RenderDrawData(ImGui::GetDrawData());

        glfwSwapBuffers(window);
    }

    // 5) Temizlik
    ImGui_ImplOpenGL3_Shutdown();
    ImGui_ImplGlfw_Shutdown();
    ImGui::DestroyContext();
    glfwDestroyWindow(window);
    glfwTerminate();
    return 0;
}
\end{lstlisting}

\gsub{Satır Satır Ne Oldu?}

\begin{center}
\begin{tabularx}{\textwidth}{@{}>{\bfseries\color{anaMavi}}l >{\raggedright\arraybackslash}X@{}}
\toprule
\rowcolor{acikMavi}\color{anaMavi}\textbf{Bölüm} & \color{anaMavi}\textbf{Açıklama}\\
\midrule
CreateContext & ImGui'nin tüm iç durumunu tutan tekil (singleton) bağlamı yaratır. Her şeyden önce çağrılmalı.\\
\rowcolor{acikGri} GetIO & Girdi/çıktı yapısı: fare, klavye, ekran boyutu, DeltaTime, framerate buradadır.\\
Init...ForOpenGL & GLFW backend'i pencereye kancalar; \code{true} argümanı ImGui'nin girdi callback'lerini otomatik kurmasını sağlar.\\
\rowcolor{acikGri} Begin/End & Bir ImGui penceresi açar ve kapatır. İçindeki widget'lar bu pencereye çizilir.\\
Render + RenderDrawData & UI'yı üçgen listesine çevirir ve OpenGL ile çizer.\\
\bottomrule
\end{tabularx}
\end{center}

\uyari{\code{ImGui::Begin} her zaman bir \code{ImGui::End} ile eşleşmelidir ---
\code{Begin} \code{false} dönse bile! Bu, çoğu \code{BeginXXX}/\code{EndXXX}
çiftinden farklıdır (birçoğunda \code{Begin} false dönerse \code{End}
çağrılmaz). Pencereler için kural: \code{Begin} ne dönerse dönsün \code{End}
çağır. Güvenli kalıp: \code{if (ImGui::Begin("X")) \{ ... \} ImGui::End();}}

\gsub{Ana Döngüyü Yeniden Kullanılabilir Yapmak}

Her projede yukarıdaki iskeleti tekrar yazmak istemeyiz. Kısım 4'te bu döngüyü
bir \code{App} sınıfına saracağız; böylece sadece ``bu karede ne çizilsin''
sorusuna odaklanabileceğiz. Şimdilik yukarıdaki \code{main.cpp} tüm örnekler
için yeterli --- sadece \emph{UI kodu} bölümünü değiştireceğiz.

% #####################################################################
\gpart{2}{Widget'lar ve Düzen}
% #####################################################################

% =====================================================================
\gsec{Pencereler ve Pencere Bayrakları}
% =====================================================================

Her ImGui widget'ı bir pencerenin içinde yaşar. \code{ImGui::Begin("Ad")} bir
pencere açar; aynı ada sahip \code{Begin} çağrıları \emph{aynı} pencereye yazar.
Pencerenin başlığı aynı zamanda onun kimliğidir (bkz. ID sistemi).

\gsub{Temel Pencere Kullanımı}

\begin{lstlisting}
bool show_panel = true;

// 2. argümanla bir kapatma dugmesi (X) ekleyebiliriz:
if (ImGui::Begin("Ayarlar", &show_panel)) {
    ImGui::Text("Panel icerigi burada.");
}
ImGui::End();
// show_panel, kullanici X'e basinca false olur -> bir daha cizmezsiniz.
\end{lstlisting}

\gsub{Sık Kullanılan Pencere Bayrakları}

\code{Begin}'in üçüncü argümanı \code{ImGuiWindowFlags} bit maskesidir.

\begin{center}
\begin{tabularx}{\textwidth}{@{}>{\ttfamily\footnotesize\color{koyuGri}}l >{\raggedright\arraybackslash}X@{}}
\toprule
\rowcolor{acikMavi}\normalfont\color{anaMavi}\textbf{Bayrak} & \color{anaMavi}\textbf{Etki}\\
\midrule
NoTitleBar & Başlık çubuğunu gizler\\
\rowcolor{acikGri} NoResize & Yeniden boyutlandırmayı kapatır\\
NoMove & Pencere taşınamaz\\
\rowcolor{acikGri} NoCollapse & Katlama (çift tık) okunu kaldırır\\
AlwaysAutoResize & Pencere içeriğe göre boyutlanır\\
\rowcolor{acikGri} NoBackground & Arka planı şeffaf yapar (HUD için ideal)\\
MenuBar & Pencereye bir menü çubuğu ekler\\
\rowcolor{acikGri} NoScrollbar & Kaydırma çubuğunu gizler\\
NoSavedSettings & Konum/boyutu \code{imgui.ini}'ye kaydetmez\\
\bottomrule
\end{tabularx}
\end{center}

\begin{lstlisting}
ImGuiWindowFlags flags = ImGuiWindowFlags_NoResize
                       | ImGuiWindowFlags_NoCollapse
                       | ImGuiWindowFlags_AlwaysAutoResize;
ImGui::Begin("Sabit Panel", nullptr, flags);
ImGui::Text("Bu pencere boyutlandirilamaz ve katlanamaz.");
ImGui::End();
\end{lstlisting}

\gsub{Konum ve Boyut Ayarlama}

Pencere konumu/boyutu \code{imgui.ini}'de saklanır ve kullanıcı taşıyabilir.
Bir başlangıç değeri vermek için \code{ImGuiCond} kullanırız: \code{\_Once}
sadece ilk açılışta, \code{\_Always} her karede uygular.

\begin{lstlisting}
ImGui::SetNextWindowPos(ImVec2(20, 20), ImGuiCond_FirstUseEver);
ImGui::SetNextWindowSize(ImVec2(400, 300), ImGuiCond_FirstUseEver);
ImGui::Begin("Konumlu Pencere");
ImGui::End();
\end{lstlisting}

\uyari{Konumu \code{ImGuiCond\_Always} ile her karede zorlarsanız kullanıcı
pencereyi taşıyamaz --- taşısa da bir sonraki karede geri sıçrar. Başlangıç
konumu için \emph{daima} \code{\_FirstUseEver} ya da \code{\_Once} kullanın.}

% =====================================================================
\gsec{Temel Girdi Widget'ları}
% =====================================================================

Bu bölümde en çok kullanılan girdi widget'larını tek tek görelim. Ortak
desen: widget bir \code{bool} döner (etkileşim oldu mu?) ve bir referans
aldığı değişkeni \emph{yerinde} günceller.

\gsub{Metin ve Etiketler}

\begin{lstlisting}
ImGui::Text("Duz metin");
ImGui::TextColored(ImVec4(1,0.6f,0,1), "Turuncu vurgulu metin");
ImGui::TextDisabled("Soluk (devre disi gorunumlu) metin");
ImGui::TextWrapped("Cok uzun bir metin otomatik satir kaydirilir...");
ImGui::BulletText("Madde imli satir");
ImGui::LabelText("FPS", "%.1f", 59.9f);   // sag: etiket, sol: deger
ImGui::SeparatorText("Bolum Basligi");     // ortasinda yazi olan cizgi

// printf tarzi bicimlendirme:
int score = 42;
ImGui::Text("Skor: %d", score);
\end{lstlisting}

\gsub{Butonlar}

\begin{lstlisting}
if (ImGui::Button("Kaydet")) {
    // tiklandi
}
ImGui::SameLine();                       // sonraki widget ayni satirda
if (ImGui::Button("Iptal")) { }

ImGui::SmallButton("Kucuk");             // satir ici kucuk dugme
ImGui::Button("Genis", ImVec2(-FLT_MIN, 0)); // -FLT_MIN: mevcut genisligi doldur
ImGui::ArrowButton("ok", ImGuiDir_Right); // yon oklu dugme

// Basili tutuldukca tekrar tetiklenen dugme:
ImGui::PushButtonRepeat(true);
if (ImGui::Button("++")) counter++;
ImGui::PopButtonRepeat();
\end{lstlisting}

\gsub{Checkbox ve Radio}

\begin{lstlisting}
static bool vsync = true;
ImGui::Checkbox("VSync acik", &vsync);

static int mode = 0;
ImGui::RadioButton("Dusuk", &mode, 0); ImGui::SameLine();
ImGui::RadioButton("Orta",  &mode, 1); ImGui::SameLine();
ImGui::RadioButton("Yuksek", &mode, 2);

// Bit bayraklari icin CheckboxFlags:
static int flags = 0;
ImGui::CheckboxFlags("Bayrak A", &flags, 1 << 0);
ImGui::CheckboxFlags("Bayrak B", &flags, 1 << 1);
\end{lstlisting}

\notk[static neden?]{Örneklerde \code{static} kullanmamız, değerin kareler
arası yaşaması içindir --- immediate mode'da yerel değişken her kare sıfırlanır.
\textbf{Gerçek uygulamada} bu değerleri \code{static} değil, kendi
uygulama/durum sınıfınızın üyeleri yapın. \code{static}, sadece tek örnekli
hızlı denemeler için uygundur.}

\gsub{Slider ve Drag}

Slider bir aralıkta değer seçtirir; Drag ise sınırsız (isteğe bağlı sınırlı)
değeri fare sürükleyerek değiştirir.

\begin{lstlisting}
static float volume = 0.5f;
ImGui::SliderFloat("Ses", &volume, 0.0f, 1.0f);

static int level = 3;
ImGui::SliderInt("Seviye", &level, 1, 10);

static float pos[3] = {0,0,0};
ImGui::SliderFloat3("Konum", pos, -10.0f, 10.0f);  // 3 bilesenli

static float speed = 1.0f;
ImGui::DragFloat("Hiz", &speed, 0.01f /*adim*/, 0.0f, 100.0f);

static int range[2] = {10, 90};
ImGui::DragIntRange2("Aralik", &range[0], &range[1], 1, 0, 100);

// Bicim ve logaritmik olcek:
ImGui::SliderFloat("Frekans", &speed, 1.0f, 20000.0f,
                   "%.0f Hz", ImGuiSliderFlags_Logarithmic);
\end{lstlisting}


\gsub{Metin Girişi (InputText)}

InputText, C-tarzı bir \code{char} tamponuna yazar. \code{std::string} ile
kullanmak için \code{imgui\_stdlib.h} (misc/cpp klasöründe) eklenir.

\begin{lstlisting}
static char name[128] = "";
ImGui::InputText("Isim", name, IM_ARRAYSIZE(name));

static char password[64] = "";
ImGui::InputText("Sifre", password, IM_ARRAYSIZE(password),
                 ImGuiInputTextFlags_Password);

static char notes[1024] = "";
ImGui::InputTextMultiline("Notlar", notes, IM_ARRAYSIZE(notes),
                          ImVec2(-FLT_MIN, 100));

// Sayisal giris:
static int age = 25;
ImGui::InputInt("Yas", &age);
static float pi = 3.14f;
ImGui::InputFloat("Pi", &pi, 0.01f, 1.0f, "%.3f");
\end{lstlisting}

\ipucu{\code{std::string} kullanmak için: \code{\#include "misc/cpp/imgui\_stdlib.h"}
ekleyin ve \code{std::string s; ImGui::InputText("X", \&s);} yazın. Bu, tampon
boyutu derdinden kurtarır. Kısım 3'te InputText callback'lerini de göreceğiz.}

\gsub{Combo (Açılır Liste) ve ListBox}

\begin{lstlisting}
// Basit yol: null ile ayrilmis tek string
const char* items = "Kirmizi\0Yesil\0Mavi\0\0";
static int current = 0;
ImGui::Combo("Renk", &current, items);

// Dizi ile:
const char* names[] = {"Dusuk", "Orta", "Yuksek", "Ultra"};
static int quality = 1;
ImGui::Combo("Kalite", &quality, names, IM_ARRAYSIZE(names));

// Tam kontrol icin BeginCombo/EndCombo:
static int sel = 0;
if (ImGui::BeginCombo("Gelismis", names[sel])) {
    for (int i = 0; i < IM_ARRAYSIZE(names); i++) {
        bool is_selected = (sel == i);
        if (ImGui::Selectable(names[i], is_selected)) sel = i;
        if (is_selected) ImGui::SetItemDefaultFocus();
    }
    ImGui::EndCombo();
}

// ListBox: birden fazla ogeyi ayni anda gorunur kutuda gosterir
static int lb = 0;
ImGui::ListBox("Liste", &lb, names, IM_ARRAYSIZE(names), 4 /*gorunur satir*/);
\end{lstlisting}

\gsub{Renk Seçiciler}

\begin{lstlisting}
static float color3[3] = {1.0f, 0.5f, 0.2f};
ImGui::ColorEdit3("Renk (RGB)", color3);

static float color4[4] = {1.0f, 0.5f, 0.2f, 1.0f};
ImGui::ColorEdit4("Renk (RGBA)", color4);

// Buyuk tekerlekli secici:
ImGui::ColorPicker4("Secici", color4);

// Sadece renk kutusu (secici acmadan):
ImGui::ColorButton("onizleme", ImVec4(color4[0],color4[1],color4[2],color4[3]));
\end{lstlisting}

% =====================================================================
\gsec{Düzen (Layout) ve Gruplama}
% =====================================================================

ImGui varsayılan olarak widget'ları alt alta dizer. Yatay düzen, boşluk,
girinti ve gruplama için özel komutlar vardır.

\gsub{Yatay Düzen ve Boşluk}

\begin{lstlisting}
ImGui::Text("A"); ImGui::SameLine(); ImGui::Text("B");  // yan yana
ImGui::SameLine(0, 40);       // aralarinda 40 piksel bosluk
ImGui::Text("C (40px sonra)");

ImGui::Spacing();             // dikey bosluk
ImGui::Separator();           // yatay ayrac cizgi
ImGui::NewLine();             // bos satir

ImGui::Indent();              // girinti ekle
ImGui::Text("Girintili");
ImGui::Unindent();            // girintiyi geri al

// Bir sonraki widget'in konumunu elle ayarla:
ImGui::SetCursorPosX(200);
ImGui::Text("x=200'de");
\end{lstlisting}

\gsub{Gruplar ve Dikey Hizalama}

\code{BeginGroup}/\code{EndGroup} bir grup widget'ı tek bir öğe gibi ele alır;
böylece \code{SameLine} onları bütün olarak yan yana koyabilir.

\begin{lstlisting}
ImGui::BeginGroup();
ImGui::Text("Sol sutun");
ImGui::Button("Dugme 1");
ImGui::Button("Dugme 2");
ImGui::EndGroup();

ImGui::SameLine();

ImGui::BeginGroup();
ImGui::Text("Sag sutun");
ImGui::Button("Dugme 3");
ImGui::EndGroup();
\end{lstlisting}

\gsub{Sütunlar: Table API}

Eski \code{Columns()} API'si yerine artık esnek \textbf{Table API}'si
önerilir. Basit bir iki sütunlu düzen:

\begin{lstlisting}
if (ImGui::BeginTable("layout", 2)) {
    ImGui::TableNextColumn(); ImGui::Text("Etiket 1");
    ImGui::TableNextColumn(); ImGui::Text("Deger 1");
    ImGui::TableNextColumn(); ImGui::Text("Etiket 2");
    ImGui::TableNextColumn(); ImGui::Text("Deger 2");
    ImGui::EndTable();
}
\end{lstlisting}

\gsub{Genişlik Kontrolü ve Öğe Genişliği}

\begin{lstlisting}
ImGui::PushItemWidth(120);        // sonraki widget'lar 120 px genis
ImGui::SliderFloat("A", &a, 0, 1);
ImGui::SliderFloat("B", &b, 0, 1);
ImGui::PopItemWidth();

ImGui::SetNextItemWidth(-100);    // saga 100 px bosluk birakip doldur
ImGui::InputText("Ara", buf, 64);
\end{lstlisting}

\bilgi{Genişlik değerlerinde işaret önemlidir: pozitif değer sabit piksel
genişliğidir; negatif değer ``sağdan şu kadar boşluk bırak, gerisini doldur''
demektir. \code{-FLT\_MIN} ise ``tüm kalan genişliği doldur'' için yaygın bir
kısayoldur.}

% =====================================================================
\gsec{ID Sistemi: ImGui'nin En Kritik Kavramı}
% =====================================================================

ImGui her widget'ı bir \textbf{ID} ile takip eder (hangi öğe aktif, hangisine
tıklandı vb.). ID varsayılan olarak widget'ın \emph{etiketinden} üretilir.
Bu, yeni başlayanların en sık düştüğü tuzağın kaynağıdır.

\gsub{Problem: Aynı Etiket = Aynı ID}

\begin{lstlisting}
// HATALI: iki dugme ayni "Sil" etiketine sahip -> ayni ID!
// ImGui ikisini tek dugme sanar; ikincisi calismayabilir.
for (int i = 0; i < 3; i++) {
    if (ImGui::Button("Sil")) { erase(i); }   // <- ID cakismasi
}
\end{lstlisting}

\gsub{Çözüm 1: \code{\#\#} ile Gizli ID}

Etikette \code{\#\#} sonrası kısım ekranda \emph{görünmez} ama ID'ye katılır.

\begin{lstlisting}
for (int i = 0; i < 3; i++) {
    // Ekranda "Sil" gorunur, ID benzersizdir:
    char label[32];
    snprintf(label, sizeof(label), "Sil##%d", i);
    if (ImGui::Button(label)) { erase(i); }
}
\end{lstlisting}

\gsub{Çözüm 2: \code{PushID}/\code{PopID}}

Daha temiz yöntem: döngü gövdesini bir ID kapsamına almak.

\begin{lstlisting}
for (int i = 0; i < 3; i++) {
    ImGui::PushID(i);              // bu iterasyona ozgu ID kapsami
    if (ImGui::Button("Sil")) { erase(i); }   // etiket ayni ama ID benzersiz
    ImGui::SameLine();
    ImGui::Text("Oge %d", i);
    ImGui::PopID();
}
\end{lstlisting}

\notk[\code{\#\#} ve \code{\#\#\#} farkı]{\code{"Baslik\#\#gizli"}: görünen
metin ``Baslik'', ID ise tüm dizeden üretilir. \code{"\#\#\#sabitID"}: görünen
metin boşalır ve \emph{sadece} \code{\#\#\#} sonrası ID belirler --- bu, etiketi
her karede değişen (örn. sayaç gösteren) bir pencerenin kimliğini sabit tutmak
için kullanılır: \code{ImGui::Begin("Skor: 42\#\#\#skorPenceresi");}}

\uyari{ID çakışması sessiz hatalara yol açar: yanlış widget aktifleşir,
tıklamalar kaybolur, InputText odağı atlar. Bir liste/döngü içinde widget
üretiyorsanız \textbf{daima} \code{PushID}/\code{PopID} kullanın. Bu, ImGui'de
en sık karşılaşılan hata sınıfıdır.}


% =====================================================================
\gsec{Kapsayıcı Widget'lar}
% =====================================================================

Bu widget'lar başka widget'ları içinde barındırır: ağaçlar, seçilebilir
listeler, menüler, sekmeler ve açılır kutular.

\gsub{Ağaç Düğümleri (TreeNode)}

\begin{lstlisting}
if (ImGui::TreeNode("Kok")) {
    ImGui::Text("Kok altindaki icerik");
    if (ImGui::TreeNode("Alt dal")) {
        ImGui::BulletText("Yaprak 1");
        ImGui::BulletText("Yaprak 2");
        ImGui::TreePop();          // her acilan TreeNode icin bir TreePop
    }
    ImGui::TreePop();
}

// Bayraklarla: varsayilan acik, tiklamayla secim
ImGuiTreeNodeFlags flags = ImGuiTreeNodeFlags_DefaultOpen
                         | ImGuiTreeNodeFlags_OpenOnArrow;
if (ImGui::TreeNodeEx("Gelismis Dal", flags)) {
    ImGui::Text("...");
    ImGui::TreePop();
}
\end{lstlisting}

\gsub{Seçilebilir Öğeler (Selectable)}

\begin{lstlisting}
static int selected = -1;
const char* files[] = {"main.cpp", "app.h", "utils.cpp", "readme.md"};
for (int i = 0; i < IM_ARRAYSIZE(files); i++) {
    if (ImGui::Selectable(files[i], selected == i)) {
        selected = i;              // tiklaninca sec
    }
    // Cift tik algila:
    if (ImGui::IsItemHovered() && ImGui::IsMouseDoubleClicked(0)) {
        open_file(files[i]);
    }
}
\end{lstlisting}

\gsub{CollapsingHeader}

Katlanabilir başlık; ayar panelleri için idealdir. \code{TreePop} gerektirmez.

\begin{lstlisting}
if (ImGui::CollapsingHeader("Grafik Ayarlari")) {
    ImGui::Checkbox("VSync", &vsync);
    ImGui::SliderInt("MSAA", &msaa, 0, 8);
}
if (ImGui::CollapsingHeader("Ses Ayarlari",
                            ImGuiTreeNodeFlags_DefaultOpen)) {
    ImGui::SliderFloat("Ana Ses", &master, 0, 1);
}
\end{lstlisting}

\gsub{Menüler ve Menü Çubuğu}

Pencere menüsü için pencereye \code{ImGuiWindowFlags\_MenuBar}, uygulama
geneli için \code{BeginMainMenuBar} kullanılır.

\begin{lstlisting}
// Uygulama ust menusu (tum ekranin ustunde):
if (ImGui::BeginMainMenuBar()) {
    if (ImGui::BeginMenu("Dosya")) {
        if (ImGui::MenuItem("Yeni", "Ctrl+N"))  { new_file(); }
        if (ImGui::MenuItem("Ac",   "Ctrl+O"))  { open(); }
        ImGui::Separator();
        if (ImGui::MenuItem("Cikis", "Alt+F4")) { quit(); }
        ImGui::EndMenu();
    }
    if (ImGui::BeginMenu("Duzen")) {
        static bool wrap = true;
        ImGui::MenuItem("Satir kaydir", nullptr, &wrap); // isaretli oge
        ImGui::EndMenu();
    }
    ImGui::EndMainMenuBar();
}
\end{lstlisting}

\gsub{Sekmeler (Tab Bar)}

\begin{lstlisting}
if (ImGui::BeginTabBar("Sekmeler")) {
    if (ImGui::BeginTabItem("Genel")) {
        ImGui::Text("Genel ayarlar");
        ImGui::EndTabItem();
    }
    if (ImGui::BeginTabItem("Ag")) {
        ImGui::Text("Ag ayarlari");
        ImGui::EndTabItem();
    }
    if (ImGui::BeginTabItem("Hakkinda")) {
        ImGui::Text("Surum 1.0");
        ImGui::EndTabItem();
    }
    ImGui::EndTabBar();
}
\end{lstlisting}

\gsub{Açılır Pencereler ve Modal'lar (Popup)}

\begin{lstlisting}
// Sag tik menusu:
if (ImGui::Button("Sag-tik menusu ac")) ImGui::OpenPopup("baglam");
if (ImGui::BeginPopup("baglam")) {
    if (ImGui::MenuItem("Kopyala")) {}
    if (ImGui::MenuItem("Yapistir")) {}
    ImGui::EndPopup();
}

// Modal onay penceresi:
if (ImGui::Button("Sil")) ImGui::OpenPopup("Emin misiniz?");
if (ImGui::BeginPopupModal("Emin misiniz?", nullptr,
                           ImGuiWindowFlags_AlwaysAutoResize)) {
    ImGui::Text("Bu islem geri alinamaz.");
    if (ImGui::Button("Evet", ImVec2(120,0))) {
        do_delete();
        ImGui::CloseCurrentPopup();
    }
    ImGui::SameLine();
    if (ImGui::Button("Hayir", ImVec2(120,0))) ImGui::CloseCurrentPopup();
    ImGui::EndPopup();
}
\end{lstlisting}

\gsub{İpucu Balonları (Tooltip)}

\begin{lstlisting}
ImGui::Text("Uzerine gel (?)");
if (ImGui::IsItemHovered()) {
    ImGui::SetTooltip("Bu bir yardim balonudur.");
}
// Zengin icerikli tooltip:
ImGui::Button("Detay");
if (ImGui::IsItemHovered()) {
    ImGui::BeginTooltip();
    ImGui::Text("Baslik");
    ImGui::Separator();
    ImGui::BulletText("Satir 1");
    ImGui::BulletText("Satir 2");
    ImGui::EndTooltip();
}
\end{lstlisting}

% =====================================================================
\gsec{Tablolar: Güçlü Veri Gösterimi}
% =====================================================================

Table API'si ImGui'nin en yeteneklidir: sıralama, yeniden boyutlandırma,
dondurma (freeze), sütun gizleme ve kaydırma destekler. Elektronik tablo
benzeri araçların temelidir.

\gsub{Temel Tablo}

\begin{lstlisting}
ImGuiTableFlags flags = ImGuiTableFlags_Borders
                      | ImGuiTableFlags_RowBg
                      | ImGuiTableFlags_Resizable;
if (ImGui::BeginTable("kisiler", 3, flags)) {
    // Baslik satiri
    ImGui::TableSetupColumn("Ad");
    ImGui::TableSetupColumn("Yas");
    ImGui::TableSetupColumn("Sehir");
    ImGui::TableHeadersRow();

    struct Row { const char* ad; int yas; const char* sehir; };
    Row rows[] = {{"Ada",30,"Izmir"},{"Kerem",25,"Ankara"},{"Zeynep",28,"Bursa"}};
    for (auto& r : rows) {
        ImGui::TableNextRow();
        ImGui::TableSetColumnIndex(0); ImGui::Text("%s", r.ad);
        ImGui::TableSetColumnIndex(1); ImGui::Text("%d", r.yas);
        ImGui::TableSetColumnIndex(2); ImGui::Text("%s", r.sehir);
    }
    ImGui::EndTable();
}
\end{lstlisting}

\gsub{Sıralanabilir ve Dondurulmuş Başlıklı Tablo}

\begin{lstlisting}
ImGuiTableFlags flags = ImGuiTableFlags_Borders | ImGuiTableFlags_RowBg
                      | ImGuiTableFlags_Sortable
                      | ImGuiTableFlags_ScrollY;
if (ImGui::BeginTable("veri", 3, flags, ImVec2(0, 200))) {
    ImGui::TableSetupScrollFreeze(0, 1);   // ilk satiri dondur (kayarken sabit)
    ImGui::TableSetupColumn("ID",  ImGuiTableColumnFlags_DefaultSort);
    ImGui::TableSetupColumn("Ad");
    ImGui::TableSetupColumn("Puan");
    ImGui::TableHeadersRow();

    // ImGuiTableSortSpecs* specs = ImGui::TableGetSortSpecs();
    // -> specs->Specs[0].ColumnIndex ve SortDirection ile verinizi siralayin

    for (int i = 0; i < 100; i++) {
        ImGui::TableNextRow();
        ImGui::TableNextColumn(); ImGui::Text("%d", i);
        ImGui::TableNextColumn(); ImGui::Text("Oge %d", i);
        ImGui::TableNextColumn(); ImGui::Text("%d", (i * 37) % 100);
    }
    ImGui::EndTable();
}
\end{lstlisting}

\ipucu{Çok satırlı tablolarda (binlerce satır) her satırı çizmek yerine
\code{ImGuiListClipper} kullanın: sadece görünür satırlar çizilir. Kısım 3'te
detaylandırıyoruz. Table + Clipper birlikte, milyonlarca satırlık veriyi
akıcı gösterebilir.}

% =====================================================================
\gsec{Grafikler ve Görselleştirme}
% =====================================================================

ImGui çekirdeği basit çizgi/histogram grafikleri sunar. Daha ciddi
görselleştirme için ImPlot eklentisi kullanılır, ama çekirdek araçlar hızlı
hata ayıklama için yeterlidir.

\gsub{Çizgi ve Histogram}

\begin{lstlisting}
float samples[90];
for (int i = 0; i < 90; i++)
    samples[i] = sinf(i * 0.2f + ImGui::GetTime());

ImGui::PlotLines("Sinus", samples, IM_ARRAYSIZE(samples),
                 0, "canli", -1.0f, 1.0f, ImVec2(0, 80));

float bars[] = {0.2f, 0.7f, 0.4f, 0.9f, 0.5f, 0.3f};
ImGui::PlotHistogram("Cubuklar", bars, IM_ARRAYSIZE(bars),
                     0, nullptr, 0.0f, 1.0f, ImVec2(0, 80));
\end{lstlisting}

\gsub{İlerleme Çubuğu}

\begin{lstlisting}
static float progress = 0.0f;
progress += ImGui::GetIO().DeltaTime * 0.3f;
if (progress > 1.0f) progress = 0.0f;

ImGui::ProgressBar(progress, ImVec2(-FLT_MIN, 0));
ImGui::ProgressBar(progress, ImVec2(-FLT_MIN, 0), "Yukleniyor...");
\end{lstlisting}

\notk[ImPlot ne zaman?]{Kaydırılabilir/yakınlaştırılabilir eksenler, birden
çok seri, gerçek zamanlı kayan grafik, saçılım (scatter), ısı haritası
istiyorsanız ImPlot (\url{https://github.com/epezent/implot}) kullanın.
Çekirdeğin \code{PlotLines}'ı yalnızca ``şu diziye hızlıca bir göz atayım''
durumları içindir.}


% #####################################################################
\gpart{3}{İleri Konular}
% #####################################################################

% =====================================================================
\gsec{Stil ve Tema}
% =====================================================================

ImGui'nin görünümü \code{ImGuiStyle} yapısı üzerinden tamamen özelleştirilir.
Renkler, kenar yuvarlaklığı, boşluklar, kenarlık kalınlıkları hepsi buradadır.

\gsub{Hazır Temalar ve Genel Ayarlar}

\begin{lstlisting}
ImGui::StyleColorsDark();     // koyu (varsayilan)
ImGui::StyleColorsLight();    // acik
ImGui::StyleColorsClassic();  // eski klasik

ImGuiStyle& style = ImGui::GetStyle();
style.WindowRounding    = 6.0f;   // pencere koseleri
style.FrameRounding     = 4.0f;   // dugme/kutu koseleri
style.GrabRounding      = 4.0f;   // slider tutamaci
style.WindowPadding     = ImVec2(10, 10);
style.ItemSpacing       = ImVec2(8, 6);
style.ScrollbarSize     = 14.0f;
style.WindowBorderSize  = 1.0f;
\end{lstlisting}

\gsub{Renkleri Elle Ayarlamak}

\begin{lstlisting}
ImVec4* colors = ImGui::GetStyle().Colors;
colors[ImGuiCol_WindowBg]       = ImVec4(0.10f, 0.10f, 0.12f, 1.00f);
colors[ImGuiCol_Button]         = ImVec4(0.20f, 0.45f, 0.75f, 1.00f);
colors[ImGuiCol_ButtonHovered]  = ImVec4(0.26f, 0.55f, 0.90f, 1.00f);
colors[ImGuiCol_ButtonActive]   = ImVec4(0.15f, 0.38f, 0.65f, 1.00f);
colors[ImGuiCol_Header]         = ImVec4(0.20f, 0.45f, 0.75f, 0.55f);
colors[ImGuiCol_FrameBg]        = ImVec4(0.16f, 0.16f, 0.20f, 1.00f);
colors[ImGuiCol_Text]           = ImVec4(0.90f, 0.90f, 0.92f, 1.00f);
\end{lstlisting}

\gsub{Yerel (Geçici) Stil Değişikliği: Push/Pop}

Belirli bir widget için stili geçici değiştirmek istersek Push/Pop kullanırız.
\textbf{Her Push bir Pop ile eşleşmelidir}, yoksa stil yığını bozulur.

\begin{lstlisting}
// Sadece bu dugme kirmizi olsun:
ImGui::PushStyleColor(ImGuiCol_Button,        ImVec4(0.7f,0.15f,0.15f,1));
ImGui::PushStyleColor(ImGuiCol_ButtonHovered, ImVec4(0.85f,0.2f,0.2f,1));
if (ImGui::Button("TEHLIKELI: Sil")) { }
ImGui::PopStyleColor(2);   // 2 renk push edildi -> 2 pop

// Gecici deger degisikligi:
ImGui::PushStyleVar(ImGuiStyleVar_FrameRounding, 12.0f);
ImGui::Button("Yuvarlak dugme");
ImGui::PopStyleVar();
\end{lstlisting}

\uyari{Push/Pop sayıları tutmazsa hata ID sisteminden bile kötüdür: stil yığını
kareler arası taşar ve tüm arayüz bozulur ya da assert atar. Kural: her
\code{PushStyleColor(n renk)} için \code{PopStyleColor(n)}. Bir kod yolunda
\code{return} varsa Pop'u atlamamaya dikkat edin (RAII sarıcı yazmak iyi fikir).}

% =====================================================================
\gsec{Fontlar ve Türkçe Karakter Desteği}
% =====================================================================

ImGui'nin varsayılan fontu yalnızca temel ASCII içerir; Türkçe ``ç, ğ, ı, İ,
ş, ö, ü'' karakterleri görünmez. Bunu çözmek için bir TTF font yükleyip doğru
\emph{glyph aralığını} belirtmeliyiz.

\gsub{Türkçe Fontu Yükleme}

\begin{lstlisting}
ImGuiIO& io = ImGui::GetIO();

// Turkce dahil Latin karakterleri iceren glyph araligi:
static const ImWchar ranges[] = {
    0x0020, 0x00FF,   // Temel Latin + Latin-1 (aksanli harfler)
    0x0100, 0x017F,   // Latin Extended-A (s-cedilla, g-breve: s, g)
    0,
};
ImFontConfig cfg;
cfg.OversampleH = 2;
cfg.OversampleV = 1;

io.Fonts->AddFontFromFileTTF(
    "fonts/DejaVuSans.ttf", 18.0f, &cfg, ranges);

// Font atlasi backend baslatilmadan ONCE olusturulmali; genelde
// ImGui_ImplOpenGL3_Init otomatik yapar. Elle: io.Fonts->Build();
\end{lstlisting}

\bilgi{Türkçe için kritik karakterler: \code{ğ (U+011F)}, \code{Ğ (U+011E)},
\code{ş (U+015F)}, \code{Ş (U+015E)}, \code{ı (U+0131)}, \code{İ (U+0130)}.
Bunlar Latin Extended-A (0x0100--0x017F) aralığındadır. \code{ç, ö, ü} ise
Latin-1'de (0x00FF'e kadar) yer alır. İkisini de eklemezseniz kelimeler
eksik görünür. Kod dosyanız da UTF-8 kaydedilmeli.}

\gsub{Birden Fazla Font ve Simge Fontları}

\begin{lstlisting}
ImFont* font_default = io.Fonts->AddFontDefault();
ImFont* font_big = io.Fonts->AddFontFromFileTTF("fonts/Roboto.ttf", 32.0f);

// Kullanirken:
ImGui::PushFont(font_big);
ImGui::Text("BUYUK BASLIK");
ImGui::PopFont();

// Ikon fontlarini (FontAwesome) metin fontuyla birlestirmek:
ImFontConfig icons_cfg;
icons_cfg.MergeMode = true;   // onceki fontun uzerine birlestir
static const ImWchar icon_ranges[] = { 0xf000, 0xf3ff, 0 };
io.Fonts->AddFontFromFileTTF("fonts/fa-solid.ttf", 16.0f,
                             &icons_cfg, icon_ranges);
\end{lstlisting}

% =====================================================================
\gsec{Docking ve Çoklu Viewport}
% =====================================================================

Docking, pencereleri birbirine ve ana pencerenin kenarlarına yapıştırma
özelliğidir --- IDE benzeri düzenler yaratır. Viewport ise ImGui pencerelerini
işletim sistemi penceresi olarak ana pencerenin \emph{dışına} taşımaya izin
verir. İkisi de ImGui'nin \code{docking} dalında bulunur.

\gsub{Docking'i Etkinleştirme}

\begin{lstlisting}
ImGuiIO& io = ImGui::GetIO();
io.ConfigFlags |= ImGuiConfigFlags_DockingEnable;        // docking
io.ConfigFlags |= ImGuiConfigFlags_ViewportsEnable;      // dis pencereler

// Viewport acikken render dongusunun sonuna eklenir:
ImGui::Render();
// ... normal cizim ...
if (io.ConfigFlags & ImGuiConfigFlags_ViewportsEnable) {
    GLFWwindow* backup = glfwGetCurrentContext();
    ImGui::UpdatePlatformWindows();
    ImGui::RenderPlatformWindowsDefault();
    glfwMakeContextCurrent(backup);
}
\end{lstlisting}

\gsub{Ana DockSpace Oluşturma}

Tüm ekranı kaplayan, diğer pencerelerin yapışacağı bir dock alanı:

\begin{lstlisting}
void DrawDockSpace() {
    ImGuiViewport* vp = ImGui::GetMainViewport();
    ImGui::SetNextWindowPos(vp->WorkPos);
    ImGui::SetNextWindowSize(vp->WorkSize);
    ImGui::SetNextWindowViewport(vp->ID);

    ImGuiWindowFlags flags = ImGuiWindowFlags_NoTitleBar
        | ImGuiWindowFlags_NoResize | ImGuiWindowFlags_NoMove
        | ImGuiWindowFlags_NoBringToFrontOnFocus | ImGuiWindowFlags_NoNavFocus
        | ImGuiWindowFlags_MenuBar;

    ImGui::PushStyleVar(ImGuiStyleVar_WindowRounding, 0.0f);
    ImGui::PushStyleVar(ImGuiStyleVar_WindowPadding, ImVec2(0, 0));
    ImGui::Begin("DockSpaceRoot", nullptr, flags);
    ImGui::PopStyleVar(2);

    ImGuiID dock_id = ImGui::GetID("MainDockSpace");
    ImGui::DockSpace(dock_id, ImVec2(0, 0), ImGuiDockNodeFlags_None);

    ImGui::End();
}
\end{lstlisting}

\notk[docking dalı]{Docking ve viewport, ImGui'nin \code{docking} branch'inde
yer alır (ana \code{master}'da değil). FetchContent kullanıyorsanız
\code{GIT\_TAG}'i \code{docking} yapın veya docking özellikli bir sürüm etiketi
seçin. API'ler kararlıdır ve üretimde yaygın kullanılır.}


% =====================================================================
\gsec{DrawList ile Özel Çizim}
% =====================================================================

Widget'lar yetmediğinde \code{ImDrawList} ile doğrudan çizgi, dikdörtgen,
daire, metin ve dolgu çizebilirsiniz. Bu, özel grafikler, düğüm editörleri,
oyun mini haritaları ve tuval uygulamalarının temelidir.

\gsub{DrawList'e Erişim}

\begin{lstlisting}
// Mevcut pencerenin cizim listesi (widget'larin arkasi/onu):
ImDrawList* dl = ImGui::GetWindowDrawList();

// Tum ekranin ustune cizmek icin (overlay):
ImDrawList* fg = ImGui::GetForegroundDrawList();
ImDrawList* bg = ImGui::GetBackgroundDrawList();
\end{lstlisting}

\gsub{Temel Şekiller}

Koordinatlar \textbf{ekran uzayındadır} (pencere değil). Genelde bir taban
nokta alıp ona göre çizeriz.

\begin{lstlisting}
ImVec2 p = ImGui::GetCursorScreenPos();   // cizim baslangici (ekran uzayi)
ImU32 white = IM_COL32(255,255,255,255);
ImU32 blue  = IM_COL32( 60,140,230,255);

dl->AddLine(p, ImVec2(p.x+100, p.y+50), white, 2.0f);
dl->AddRect(ImVec2(p.x, p.y), ImVec2(p.x+80, p.y+80), white, 4.0f, 0, 1.5f);
dl->AddRectFilled(ImVec2(p.x+90,p.y), ImVec2(p.x+170,p.y+80), blue, 4.0f);
dl->AddCircle(ImVec2(p.x+220,p.y+40), 40, white, 32, 2.0f);
dl->AddCircleFilled(ImVec2(p.x+320,p.y+40), 40, blue, 32);
dl->AddTriangleFilled(ImVec2(p.x,p.y+120), ImVec2(p.x+60,p.y+120),
                      ImVec2(p.x+30,p.y+90), white);
dl->AddText(ImVec2(p.x, p.y+140), white, "DrawList metni");
dl->AddBezierCubic(p, ImVec2(p.x+50,p.y-40), ImVec2(p.x+100,p.y+80),
                   ImVec2(p.x+150,p.y), blue, 2.0f);

// Cizdigimiz alanin yer kaplamasi icin bir kukla oge birak:
ImGui::Dummy(ImVec2(360, 180));
\end{lstlisting}

\gsub{Etkileşimli Tuval Deseni}

Çizim alanını tıklanabilir yapmak için görünmez bir buton (\code{InvisibleButton})
kullanılır; sonra fare konumunu alıp çizeriz.

\begin{lstlisting}
ImVec2 canvas_p0 = ImGui::GetCursorScreenPos();
ImVec2 canvas_sz = ImGui::GetContentRegionAvail();
ImVec2 canvas_p1 = ImVec2(canvas_p0.x + canvas_sz.x,
                          canvas_p0.y + canvas_sz.y);

ImDrawList* dl = ImGui::GetWindowDrawList();
dl->AddRectFilled(canvas_p0, canvas_p1, IM_COL32(30,30,35,255));

ImGui::InvisibleButton("canvas", canvas_sz);
bool hovered = ImGui::IsItemHovered();
bool active  = ImGui::IsItemActive();

if (active) {
    ImVec2 m = ImGui::GetIO().MousePos;
    dl->AddCircleFilled(m, 5, IM_COL32(255,180,0,255));
}
\end{lstlisting}

\ipucu{DrawList'e çizerken \code{PushClipRect}/\code{PopClipRect} ile çizimi
bir dikdörtgene sınırlayabilirsiniz --- tuval dışına taşmayı engeller.
Ayrıca \code{dl->ChannelsSplit}/\code{ChannelsMerge} ile çizim katmanlarını
(örneğin düğümlerin arkasındaki bağlantı çizgileri) yönetebilirsiniz.}

% =====================================================================
\gsec{InputText Geri Çağırmaları (Callback)}
% =====================================================================

InputText'in davranışını, bir callback vererek özelleştirebiliriz: otomatik
tamamlama, geçmiş (history), karakter filtreleme veya her değişimde tetikleme.

\gsub{Callback Bayrakları ve İskelet}

\begin{lstlisting}
struct InputState { std::vector<std::string> history; int pos = -1; };

int TextCallback(ImGuiInputTextCallbackData* data) {
    InputState* st = (InputState*)data->UserData;
    switch (data->EventFlag) {
    case ImGuiInputTextFlags_CallbackCompletion:
        // TAB'a basildi: otomatik tamamlama yap
        data->InsertChars(data->CursorPos, "_tamamlandi");
        break;
    case ImGuiInputTextFlags_CallbackHistory:
        // Yukari/asagi ok: gecmiste gez
        if (data->EventKey == ImGuiKey_UpArrow && st->pos + 1 < (int)st->history.size()) {
            st->pos++;
            data->DeleteChars(0, data->BufTextLen);
            data->InsertChars(0, st->history[st->pos].c_str());
        }
        break;
    case ImGuiInputTextFlags_CallbackCharFilter:
        // Sadece rakam kabul et:
        if (data->EventChar < '0' || data->EventChar > '9') return 1; // reddet
        break;
    }
    return 0;
}

// Kullanim:
static char buf[256] = "";
static InputState state;
ImGuiInputTextFlags flags = ImGuiInputTextFlags_CallbackCompletion
                          | ImGuiInputTextFlags_CallbackHistory
                          | ImGuiInputTextFlags_EnterReturnsTrue;
if (ImGui::InputText("Komut", buf, sizeof(buf), flags, TextCallback, &state)) {
    // Enter'a basildi
    state.history.insert(state.history.begin(), buf);
    state.pos = -1;
    buf[0] = '\0';
}
\end{lstlisting}

% =====================================================================
\gsec{Klavye, Fare ve Öğe Durumu Sorgulama}
% =====================================================================

ImGui, en son çizilen öğe hakkında zengin sorgu fonksiyonları sunar. Bunlar
etkileşimli arayüzlerin kalbidir.

\gsub{Öğe Durumu Fonksiyonları}

\begin{lstlisting}
ImGui::Button("Test");
// Bir onceki widget hakkinda:
bool hovered  = ImGui::IsItemHovered();   // fare uzerinde mi
bool clicked  = ImGui::IsItemClicked();   // bu kare tiklandi mi
bool active   = ImGui::IsItemActive();    // basili tutuluyor mu
bool focused  = ImGui::IsItemFocused();   // klavye odaginda mi
bool edited   = ImGui::IsItemEdited();    // degeri bu kare degisti mi
ImVec2 rmin   = ImGui::GetItemRectMin();  // ogenin sinirlari
ImVec2 rmax   = ImGui::GetItemRectMax();

// Pencere hakkinda:
bool win_hovered = ImGui::IsWindowHovered();
bool win_focused = ImGui::IsWindowFocused();
\end{lstlisting}

\gsub{Doğrudan Klavye ve Fare}

\begin{lstlisting}
ImGuiIO& io = ImGui::GetIO();

// Fare:
ImVec2 mouse = io.MousePos;
bool left_down = ImGui::IsMouseDown(ImGuiMouseButton_Left);
bool right_click = ImGui::IsMouseClicked(ImGuiMouseButton_Right);
float wheel = io.MouseWheel;               // scroll miktari
ImVec2 drag = ImGui::GetMouseDragDelta(ImGuiMouseButton_Left);

// Klavye (yeni API):
if (ImGui::IsKeyPressed(ImGuiKey_Space)) { jump(); }
if (ImGui::IsKeyDown(ImGuiKey_W)) { move_forward(); }
if (ImGui::IsKeyChordPressed(ImGuiMod_Ctrl | ImGuiKey_S)) { save(); }
\end{lstlisting}

\uyari{Oyun/uygulama girdisiyle ImGui girdisi çakışabilir. \code{io.WantCaptureMouse}
ve \code{io.WantCaptureKeyboard} bayraklarını kontrol edin: bunlar \code{true}
iken girdi ImGui'ye aittir (örn. kullanıcı bir InputText'e yazıyor), oyununuz
o girdiyi \emph{yok saymalıdır}. Aksi hâlde metin yazarken karakter hareket eder.}

% =====================================================================
\gsec{Performans: ListClipper ve İpuçları}
% =====================================================================

Immediate mode her kare tüm UI'yı çalıştırır. Çoğu araç için bu sorunsuzdur,
ama binlerce satırlık listelerde görünmeyen öğeleri çizmek israftır.
\code{ImGuiListClipper} yalnızca görünür aralığı çizer.

\gsub{ListClipper Kullanımı}

\begin{lstlisting}
// 1.000.000 satirlik liste -- yalnizca gorunenler cizilir
ImGuiListClipper clipper;
clipper.Begin(1000000);
while (clipper.Step()) {
    for (int i = clipper.DisplayStart; i < clipper.DisplayEnd; i++) {
        ImGui::Text("Satir %d: veri...", i);
    }
}
\end{lstlisting}

\gsub{Genel Performans İpuçları}

\begin{center}
\begin{tabularx}{\textwidth}{@{}>{\bfseries\color{anaMavi}}l >{\raggedright\arraybackslash}X@{}}
\toprule
\rowcolor{acikMavi}\color{anaMavi}\textbf{İpucu} & \color{anaMavi}\textbf{Neden}\\
\midrule
Ağır işi UI'dan ayır & Dosya/ağ/hesap işini UI kodunda yapma; sonucu göster.\\
\rowcolor{acikGri} ListClipper kullan & Uzun listelerde görünmeyeni çizme.\\
Gereksiz \code{sprintf} azalt & Her karede binlerce format çağrısı birikir.\\
\rowcolor{acikGri} Kapalı pencereyi erken çık & \code{if(!Begin(...))\{End();return;\}} ile içeriği atla.\\
Clip'i backend'e bırak & ImGui zaten görünmeyen pencereyi çizmez.\\
\bottomrule
\end{tabularx}
\end{center}

\bilgi{ImGui'nin kendisi neredeyse hiç yavaşlamaz --- darboğaz genelde
\emph{sizin} UI kodunuzun her kare yaptığı gereksiz hesaplardır. Şüpheye
düşerseniz \code{ImGui::ShowMetricsWindow()} penceresini açın: kaç vertex, kaç
pencere, kaç çizim çağrısı ürettiğinizi canlı gösterir.}

\gsub{Yerleşik Demo ve Yardımcı Pencereler}

\begin{lstlisting}
static bool show_demo = true;
if (show_demo) ImGui::ShowDemoWindow(&show_demo);  // her seyi gosteren demo
ImGui::ShowMetricsWindow(nullptr);                 // performans/ic durum
ImGui::ShowStyleEditor();                          // canli stil duzenleyici
ImGui::ShowUserGuide();                            // kisayol rehberi
\end{lstlisting}

\ipucu{\code{ShowDemoWindow} bir öğrenme aracından çok daha fazlasıdır: yeni
bir widget'ı nasıl kullanacağınızı unuttuğunuzda demo penceresini açın, ilgili
bölümü bulun ve sağ üstteki menüden ``item picker'' ile o widget'ın
\code{imgui\_demo.cpp}'deki kaynak satırına gidin.}


% #####################################################################
\gpart{4}{Projeler}
% #####################################################################

Bu kısımda öğrendiğimiz her şeyi \textbf{tam, çalışan projelere} dönüştürüyoruz.
Önce tüm projelerin paylaşacağı yeniden kullanılabilir bir uygulama iskeleti
kuruyoruz; sonra her proje yalnızca bir ``çiz'' fonksiyonu yazarak bu iskelete
bağlanıyor. Projeler basitten karmaşığa doğru sıralanmıştır.

% =====================================================================
\gsec{Proje 0: Yeniden Kullanılabilir Uygulama İskeleti}
% =====================================================================

Kısım 1'deki \code{main.cpp}'yi her seferinde tekrar yazmamak için pencere
kurulumunu, ana döngüyü ve temizliği bir sınıfa saralım. Bundan sonraki tüm
projeler sadece \code{OnFrame()} fonksiyonunu doldurur.

\gsub{app.h --- İskelet}

\begin{lstlisting}
// app.h -- tum projelerin ortak temeli
#pragma once
#include "imgui.h"
#include "imgui_impl_glfw.h"
#include "imgui_impl_opengl3.h"
#include <GLFW/glfw3.h>
#include <functional>
#include <cstdio>

class App {
public:
    App(const char* title, int w = 1280, int h = 720) {
        glfwSetErrorCallback([](int e, const char* d){
            std::fprintf(stderr, "GLFW %d: %s\n", e, d); });
        glfwInit();
        glfwWindowHint(GLFW_CONTEXT_VERSION_MAJOR, 3);
        glfwWindowHint(GLFW_CONTEXT_VERSION_MINOR, 3);
        glfwWindowHint(GLFW_OPENGL_PROFILE, GLFW_OPENGL_CORE_PROFILE);
        window_ = glfwCreateWindow(w, h, title, nullptr, nullptr);
        glfwMakeContextCurrent(window_);
        glfwSwapInterval(1);

        IMGUI_CHECKVERSION();
        ImGui::CreateContext();
        ImGui::GetIO().ConfigFlags |= ImGuiConfigFlags_NavEnableKeyboard;
        ImGui::StyleColorsDark();
        ImGui_ImplGlfw_InitForOpenGL(window_, true);
        ImGui_ImplOpenGL3_Init("#version 330");
    }

    ~App() {
        ImGui_ImplOpenGL3_Shutdown();
        ImGui_ImplGlfw_Shutdown();
        ImGui::DestroyContext();
        glfwDestroyWindow(window_);
        glfwTerminate();
    }

    // Her kare cagirilacak cizim fonksiyonunu alir ve dongusu surer.
    void Run(const std::function<void()>& on_frame) {
        while (!glfwWindowShouldClose(window_)) {
            glfwPollEvents();
            ImGui_ImplOpenGL3_NewFrame();
            ImGui_ImplGlfw_NewFrame();
            ImGui::NewFrame();

            on_frame();                       // <-- projenin UI kodu

            ImGui::Render();
            int dw, dh;
            glfwGetFramebufferSize(window_, &dw, &dh);
            glViewport(0, 0, dw, dh);
            glClearColor(0.09f, 0.10f, 0.12f, 1.0f);
            glClear(GL_COLOR_BUFFER_BIT);
            ImGui_ImplOpenGL3_RenderDrawData(ImGui::GetDrawData());
            glfwSwapBuffers(window_);
        }
    }

private:
    GLFWwindow* window_ = nullptr;
};
\end{lstlisting}

\gsub{Kullanım Kalıbı}

Bundan sonra her proje bu iki satırla başlar; biz sadece köşeli parantezin
içini doldururuz:

\begin{lstlisting}
#include "app.h"
int main() {
    App app("Proje Adi");
    app.Run([&]() {
        // ... bu projenin her kare cizilen UI kodu ...
    });
}
\end{lstlisting}

% =====================================================================
\gsec{Proje 1: Yapılacaklar (To-Do) Uygulaması}
% =====================================================================

İlk projemiz klasik bir ``yapılacaklar listesi''. Ekleme, tamamlama olarak
işaretleme, silme ve filtreleme yapar. Immediate mode'un veri-yansıtma
felsefesini net gösterir: UI, \code{tasks} vektörünün her karede çizilen
görüntüsüdür.

\gsub{Tam Kaynak}

\begin{lstlisting}
#include "app.h"
#include <vector>
#include <string>

struct Task {
    std::string text;
    bool done = false;
};

int main() {
    App app("Yapilacaklar");
    std::vector<Task> tasks;
    char input[256] = "";
    int filter = 0;   // 0: tumu, 1: aktif, 2: tamamlanan

    app.Run([&]() {
        ImGui::SetNextWindowPos(ImVec2(0,0), ImGuiCond_Once);
        ImGui::SetNextWindowSize(ImGui::GetMainViewport()->Size, ImGuiCond_Once);
        ImGui::Begin("Gorevlerim", nullptr, ImGuiWindowFlags_NoResize
                     | ImGuiWindowFlags_NoMove | ImGuiWindowFlags_NoCollapse);

        // --- Ekleme satiri ---
        ImGui::SetNextItemWidth(-100);
        bool enter = ImGui::InputText("##yeni", input, sizeof(input),
                                      ImGuiInputTextFlags_EnterReturnsTrue);
        ImGui::SameLine();
        bool add = ImGui::Button("Ekle", ImVec2(-FLT_MIN, 0));
        if ((enter || add) && input[0] != '\0') {
            tasks.push_back({input, false});
            input[0] = '\0';
        }

        // --- Filtre ---
        ImGui::RadioButton("Tumu", &filter, 0);   ImGui::SameLine();
        ImGui::RadioButton("Aktif", &filter, 1);  ImGui::SameLine();
        ImGui::RadioButton("Biten", &filter, 2);
        ImGui::Separator();

        // --- Liste ---
        int to_delete = -1;
        for (int i = 0; i < (int)tasks.size(); i++) {
            Task& t = tasks[i];
            if (filter == 1 && t.done) continue;
            if (filter == 2 && !t.done) continue;

            ImGui::PushID(i);                      // benzersiz ID!
            ImGui::Checkbox("##done", &t.done);
            ImGui::SameLine();
            if (t.done) {
                ImGui::PushStyleColor(ImGuiCol_Text, ImVec4(0.5f,0.5f,0.5f,1));
                ImGui::Text("%s", t.text.c_str());
                ImGui::PopStyleColor();
            } else {
                ImGui::Text("%s", t.text.c_str());
            }
            ImGui::SameLine(ImGui::GetWindowWidth() - 50);
            if (ImGui::SmallButton("Sil")) to_delete = i;
            ImGui::PopID();
        }
        if (to_delete >= 0) tasks.erase(tasks.begin() + to_delete);

        // --- Ozet ---
        int done = 0; for (auto& t : tasks) if (t.done) done++;
        ImGui::Separator();
        ImGui::Text("Toplam: %d  |  Biten: %d  |  Kalan: %d",
                    (int)tasks.size(), done, (int)tasks.size() - done);

        ImGui::End();
    });
}
\end{lstlisting}

\notk[Neyi öğrendik?]{Silmeyi döngü \emph{içinde} yapmayıp \code{to\_delete}
değişkenine erteledik --- çünkü çizim sürerken vektörü değiştirmek indeksleri
bozar. Bu, immediate mode listelerinde standart bir kalıptır: ``döngüde niyeti
kaydet, döngüden sonra uygula''. Ayrıca her satırda \code{PushID(i)} ile ID
çakışmasını önledik.}


% =====================================================================
\gsec{Proje 2: Renkli Log Konsolu}
% =====================================================================

Hemen her araçta bir log/konsol penceresi gerekir. Bu proje; seviyeye göre
renklendirme, metin filtresi, otomatik aşağı kaydırma ve komut girişi olan
yeniden kullanılabilir bir konsol sınıfı sunar.

\gsub{Tam Kaynak}

\begin{lstlisting}
#include "app.h"
#include <vector>
#include <string>
#include <cstdarg>

enum LogLevel { LOG_INFO, LOG_WARN, LOG_ERROR };

struct Console {
    struct Line { LogLevel level; std::string text; };
    std::vector<Line> lines;
    char input[256] = "";
    ImGuiTextFilter filter;
    bool auto_scroll = true;
    bool show_info = true, show_warn = true, show_error = true;

    void Add(LogLevel lv, const char* fmt, ...) {
        char buf[1024];
        va_list args; va_start(args, fmt);
        vsnprintf(buf, sizeof(buf), fmt, args);
        va_end(args);
        lines.push_back({lv, buf});
    }

    void Draw(const char* title) {
        ImGui::Begin(title);

        // Ust arac cubugu
        if (ImGui::Button("Temizle")) lines.clear();
        ImGui::SameLine(); ImGui::Checkbox("Info", &show_info);
        ImGui::SameLine(); ImGui::Checkbox("Uyari", &show_warn);
        ImGui::SameLine(); ImGui::Checkbox("Hata", &show_error);
        ImGui::SameLine(); ImGui::Checkbox("Oto-kaydir", &auto_scroll);
        filter.Draw("Filtre", 180);
        ImGui::Separator();

        // Kaydirilabilir log alani (giris satirina yer birak)
        float footer = ImGui::GetFrameHeightWithSpacing();
        ImGui::BeginChild("scroll", ImVec2(0, -footer), false,
                          ImGuiWindowFlags_HorizontalScrollbar);
        for (auto& ln : lines) {
            if (ln.level == LOG_INFO  && !show_info)  continue;
            if (ln.level == LOG_WARN  && !show_warn)  continue;
            if (ln.level == LOG_ERROR && !show_error) continue;
            if (!filter.PassFilter(ln.text.c_str())) continue;

            ImVec4 col(0.85f,0.85f,0.85f,1);
            if (ln.level == LOG_WARN)  col = ImVec4(1.0f,0.8f,0.2f,1);
            if (ln.level == LOG_ERROR) col = ImVec4(1.0f,0.4f,0.4f,1);
            ImGui::PushStyleColor(ImGuiCol_Text, col);
            ImGui::TextUnformatted(ln.text.c_str());
            ImGui::PopStyleColor();
        }
        if (auto_scroll && ImGui::GetScrollY() >= ImGui::GetScrollMaxY())
            ImGui::SetScrollHereY(1.0f);
        ImGui::EndChild();

        ImGui::Separator();
        // Komut girisi
        ImGui::SetNextItemWidth(-60);
        bool run = ImGui::InputText("##cmd", input, sizeof(input),
                                    ImGuiInputTextFlags_EnterReturnsTrue);
        ImGui::SameLine();
        run |= ImGui::Button("Gonder");
        if (run && input[0]) {
            Add(LOG_INFO, "> %s", input);
            input[0] = '\0';
            ImGui::SetKeyboardFocusHere(-1);   // odagi girise geri ver
        }
        ImGui::End();
    }
};

int main() {
    App app("Log Konsolu");
    Console console;
    console.Add(LOG_INFO,  "Sistem baslatildi.");
    console.Add(LOG_WARN,  "Yapilandirma dosyasi bulunamadi, varsayilan kullaniliyor.");
    console.Add(LOG_ERROR, "Sunucuya baglanilamadi (port 8080).");

    float timer = 0;
    app.Run([&]() {
        // Her saniye otomatik bir log ekle (demo)
        timer += ImGui::GetIO().DeltaTime;
        if (timer > 1.0f) {
            timer = 0;
            console.Add(LOG_INFO, "Kalp atisi: t=%.1f", ImGui::GetTime());
        }
        console.Draw("Konsol");
    });
}
\end{lstlisting}

\bilgi{\code{ImGuiTextFilter}, ImGui'nin yerleşik filtre yardımcısıdır:
``\code{hata,-baglanti}'' gibi ifadeleri anlar (``hata içeren ama baglanti
içermeyen''). \code{filter.Draw()} bir arama kutusu çizer,
\code{filter.PassFilter(metin)} her satır için testi yapar. Konsollar,
listeler ve tablolarda çok işe yarar.}

% =====================================================================
\gsec{Proje 3: Canlı Sistem Monitörü}
% =====================================================================

Bu proje kayan (scrolling) grafikler çizer: sahte CPU/RAM değerlerini bir
halka tampona (ring buffer) yazıp \code{PlotLines} ile gösterir. Gerçek bir
uygulamada bu verileri işletim sisteminden okurdunuz; burada mantığa
odaklanıyoruz.

\gsub{Kayan Tampon (Scrolling Buffer)}

\begin{lstlisting}
#include "app.h"
#include <cmath>

struct ScrollBuffer {
    static const int SIZE = 120;
    float data[SIZE] = {};
    int offset = 0;

    void Push(float v) { data[offset] = v; offset = (offset + 1) % SIZE; }
    float Max() const {
        float m = 0; for (float v : data) if (v > m) m = v; return m;
    }
    float Last() const { return data[(offset - 1 + SIZE) % SIZE]; }
};

int main() {
    App app("Sistem Monitoru");
    ScrollBuffer cpu, ram;
    float t = 0;

    app.Run([&]() {
        // Sahte veri uret
        t += ImGui::GetIO().DeltaTime;
        cpu.Push(50 + 40 * sinf(t * 1.5f) + 10 * sinf(t * 7.0f));
        ram.Push(60 + 20 * sinf(t * 0.4f));

        ImGui::Begin("Kaynak Kullanimi");

        // --- CPU ---
        ImGui::Text("CPU:  %.1f%%", cpu.Last());
        ImGui::PushStyleColor(ImGuiCol_PlotLines, ImVec4(0.2f,0.8f,0.4f,1));
        ImGui::PlotLines("##cpu", cpu.data, ScrollBuffer::SIZE, cpu.offset,
                         nullptr, 0.0f, 100.0f, ImVec2(-FLT_MIN, 80));
        ImGui::PopStyleColor();

        // --- RAM ---
        ImGui::Text("RAM:  %.1f%%", ram.Last());
        ImGui::PushStyleColor(ImGuiCol_PlotLines, ImVec4(0.3f,0.6f,1.0f,1));
        ImGui::PlotLines("##ram", ram.data, ScrollBuffer::SIZE, ram.offset,
                         nullptr, 0.0f, 100.0f, ImVec2(-FLT_MIN, 80));
        ImGui::PopStyleColor();

        ImGui::Separator();
        // Ozet tablo
        if (ImGui::BeginTable("ozet", 3, ImGuiTableFlags_Borders)) {
            ImGui::TableSetupColumn("Metrik");
            ImGui::TableSetupColumn("Anlik");
            ImGui::TableSetupColumn("Tepe");
            ImGui::TableHeadersRow();
            const char* names[] = {"CPU", "RAM"};
            ScrollBuffer* bufs[] = {&cpu, &ram};
            for (int i = 0; i < 2; i++) {
                ImGui::TableNextRow();
                ImGui::TableNextColumn(); ImGui::Text("%s", names[i]);
                ImGui::TableNextColumn(); ImGui::Text("%.1f%%", bufs[i]->Last());
                ImGui::TableNextColumn(); ImGui::Text("%.1f%%", bufs[i]->Max());
            }
            ImGui::EndTable();
        }
        ImGui::Text("FPS: %.0f (%.2f ms)", ImGui::GetIO().Framerate,
                    1000.0f / ImGui::GetIO().Framerate);
        ImGui::End();
    });
}
\end{lstlisting}

\notk[Halka tampon]{\code{PlotLines}'ın \code{values\_offset} parametresi tam da
bir halka tampon için tasarlanmıştır: son yazılan konumu verirsiniz, ImGui
diziyi o noktadan itibaren dairesel okur; böylece grafik soldan sağa akar ve
en eski örnek düşer. Kaydırmalı grafiklerin en verimli yöntemidir --- her kare
diziyi kaydırmanıza gerek kalmaz.}


% =====================================================================
\gsec{Proje 4: Renk Paleti Tasarımcısı}
% =====================================================================

Renk seçicilerini, DrawList çizimini ve panoya kopyalamayı birleştiren bir
araç. Kullanıcı bir renk seçer, palete ekler, paleti kutucuklar hâlinde görür
ve seçilen rengin hex kodunu kopyalar.

\gsub{Tam Kaynak}

\begin{lstlisting}
#include "app.h"
#include <vector>

static ImU32 ToU32(const float c[4]) {
    return IM_COL32((int)(c[0]*255), (int)(c[1]*255),
                    (int)(c[2]*255), (int)(c[3]*255));
}

int main() {
    App app("Palet Tasarimcisi");
    float current[4] = {0.85f, 0.35f, 0.25f, 1.0f};
    std::vector<ImVec4> palette;

    app.Run([&]() {
        ImGui::Begin("Palet");

        // Buyuk secici
        ImGui::ColorPicker4("##picker", current,
            ImGuiColorEditFlags_NoSidePreview | ImGuiColorEditFlags_DisplayHex);

        // Hex kodu ve panoya kopyala
        int r = (int)(current[0]*255), g = (int)(current[1]*255),
            b = (int)(current[2]*255);
        char hex[16]; snprintf(hex, sizeof(hex), "#%02X%02X%02X", r, g, b);
        ImGui::Text("Hex: %s", hex);
        ImGui::SameLine();
        if (ImGui::SmallButton("Kopyala")) ImGui::SetClipboardText(hex);

        if (ImGui::Button("Palete Ekle")) {
            palette.push_back(ImVec4(current[0], current[1], current[2], 1));
        }
        ImGui::SameLine();
        if (ImGui::Button("Paleti Temizle")) palette.clear();

        ImGui::Separator();
        ImGui::Text("Palet (%d renk) -- tiklayinca sec, sag-tik sil:",
                    (int)palette.size());

        // Kutucuklari DrawList ile ciz, satir basina 8 tane
        ImVec2 start = ImGui::GetCursorScreenPos();
        ImDrawList* dl = ImGui::GetWindowDrawList();
        const float box = 36, gap = 6;
        const int per_row = 8;
        int remove = -1;
        for (int i = 0; i < (int)palette.size(); i++) {
            int col = i % per_row, row = i / per_row;
            ImVec2 p0(start.x + col*(box+gap), start.y + row*(box+gap));
            ImVec2 p1(p0.x + box, p0.y + box);

            float c[4] = {palette[i].x, palette[i].y, palette[i].z, 1};
            dl->AddRectFilled(p0, p1, ToU32(c), 4.0f);
            dl->AddRect(p0, p1, IM_COL32(255,255,255,60), 4.0f);

            // Tiklanabilir alan
            ImGui::SetCursorScreenPos(p0);
            ImGui::PushID(i);
            ImGui::InvisibleButton("sw", ImVec2(box, box));
            if (ImGui::IsItemClicked(ImGuiMouseButton_Left)) {
                current[0]=palette[i].x; current[1]=palette[i].y;
                current[2]=palette[i].z;
            }
            if (ImGui::IsItemClicked(ImGuiMouseButton_Right)) remove = i;
            ImGui::PopID();
        }
        if (remove >= 0) palette.erase(palette.begin() + remove);

        // Kutucuk alani icin yer ac
        int rows = ((int)palette.size() + per_row - 1) / per_row;
        ImGui::SetCursorScreenPos(start);
        ImGui::Dummy(ImVec2(per_row*(box+gap), rows*(box+gap) + 4));

        ImGui::End();
    });
}
\end{lstlisting}

\ipucu{\code{ImGui::SetClipboardText} ve \code{GetClipboardText} platformdan
bağımsız pano erişimi sağlar --- backend bunu işletim sistemine bağlar. Araç
geliştirirken ``hex kodu kopyala'', ``JSON'u kopyala'' gibi küçük kolaylıklar
kullanıcıyı çok mutlu eder.}

% =====================================================================
\gsec{Proje 5: Hex / Bellek Görüntüleyici}
% =====================================================================

Ters mühendislik ve dosya araçlarının klasiği: bir bayt tamponunu onaltılık
(hex) ve ASCII olarak yan yana gösteren görüntüleyici. \code{ImGuiListClipper}
sayesinde megabaytlarca veriyi akıcı gösterir --- yalnızca görünen satırlar
çizilir.

\gsub{Tam Kaynak}

\begin{lstlisting}
#include "app.h"
#include <vector>
#include <cstdint>

int main() {
    App app("Hex Goruntuleyici");

    // Ornek veri: 256 KB rastgele-benzeri desen
    std::vector<uint8_t> mem(256 * 1024);
    for (size_t i = 0; i < mem.size(); i++)
        mem[i] = (uint8_t)((i * 37 + (i >> 4)) & 0xFF);

    const int cols = 16;   // satir basina bayt

    app.Run([&]() {
        ImGui::Begin("Bellek");
        ImGui::Text("Boyut: %zu bayt (%.1f KB)", mem.size(), mem.size()/1024.0f);
        ImGui::Separator();

        int total_rows = (int)((mem.size() + cols - 1) / cols);

        ImGui::BeginChild("hexscroll", ImVec2(0,0), false,
                          ImGuiWindowFlags_HorizontalScrollbar);
        ImGui::PushFont(nullptr);  // esit genislikli font ideal olurdu

        ImGuiListClipper clipper;
        clipper.Begin(total_rows);
        while (clipper.Step()) {
            for (int row = clipper.DisplayStart; row < clipper.DisplayEnd; row++) {
                size_t base = (size_t)row * cols;

                // Adres
                ImGui::Text("%08zX: ", base);
                ImGui::SameLine();

                // Hex sutunlari
                char hexpart[64] = "";
                for (int c = 0; c < cols; c++) {
                    char b[4];
                    if (base + c < mem.size())
                        snprintf(b, sizeof(b), "%02X ", mem[base + c]);
                    else
                        snprintf(b, sizeof(b), "   ");
                    strcat(hexpart, b);
                }
                ImGui::TextUnformatted(hexpart);
                ImGui::SameLine();

                // ASCII sutunu
                char ascii[32] = " |";
                int len = 2;
                for (int c = 0; c < cols && base + c < mem.size(); c++) {
                    uint8_t ch = mem[base + c];
                    ascii[len++] = (ch >= 32 && ch < 127) ? (char)ch : '.';
                }
                ascii[len++] = '|'; ascii[len] = '\0';
                ImGui::TextUnformatted(ascii);
            }
        }
        ImGui::PopFont();
        ImGui::EndChild();
        ImGui::End();
    });
}
\end{lstlisting}

\bilgi{Buradaki püf nokta \code{ImGuiListClipper}'dır: 16.384 satırın (256 KB /
16) tamamını çizmeye çalışsaydık her kare 16 binden fazla \code{Text} çağrısı
yapardık. Clipper ile yalnızca ekrandaki $\sim$40 satır çizilir; kaydırma çubuğu
yine de tüm veriyi doğru temsil eder. Bu, ImGui'de büyük veri göstermenin
altın kuralıdır.}

\uyari{Gerçek bir hex editöründe \emph{eşit genişlikli} (monospace) font
şarttır, yoksa sütunlar hizalanmaz. Kısım 3'teki font yükleme yöntemiyle bir
monospace TTF (örn. DejaVu Sans Mono) yükleyip \code{PushFont} ile bu pencerede
kullanın.}


% =====================================================================
\gsec{Proje 6: Basit Çizim Tuvali}
% =====================================================================

DrawList ve fare etkileşimini birleştiren bir vektör çizim uygulaması.
Kullanıcı fareyle serbest çizgiler çizer (stroke), rengi ve kalınlığı seçer,
son çizgiyi geri alır ve tuvali temizler. Grafik editörlerinin çekirdek
mantığını gösterir.

\gsub{Tam Kaynak}

\begin{lstlisting}
#include "app.h"
#include <vector>

struct Stroke {
    std::vector<ImVec2> points;   // tuval-yerel koordinatlar
    ImU32 color;
    float thickness;
};

int main() {
    App app("Cizim Tuvali");
    std::vector<Stroke> strokes;
    float draw_col[4] = {1.0f, 0.8f, 0.2f, 1.0f};
    float thickness = 3.0f;
    bool drawing = false;

    app.Run([&]() {
        ImGui::Begin("Tuval");

        // Arac cubugu
        ImGui::ColorEdit3("Renk", draw_col, ImGuiColorEditFlags_NoInputs);
        ImGui::SameLine(); ImGui::SetNextItemWidth(120);
        ImGui::SliderFloat("Kalinlik", &thickness, 1.0f, 20.0f);
        ImGui::SameLine();
        if (ImGui::Button("Geri Al") && !strokes.empty()) strokes.pop_back();
        ImGui::SameLine();
        if (ImGui::Button("Temizle")) strokes.clear();

        // Tuval alani
        ImVec2 p0 = ImGui::GetCursorScreenPos();
        ImVec2 sz = ImGui::GetContentRegionAvail();
        if (sz.x < 50) sz.x = 50; if (sz.y < 50) sz.y = 50;
        ImVec2 p1(p0.x + sz.x, p0.y + sz.y);

        ImDrawList* dl = ImGui::GetWindowDrawList();
        dl->AddRectFilled(p0, p1, IM_COL32(25,25,30,255));
        dl->AddRect(p0, p1, IM_COL32(120,120,130,255));

        // Etkilesim
        ImGui::InvisibleButton("tuval", sz,
            ImGuiButtonFlags_MouseButtonLeft);
        ImVec2 mouse = ImGui::GetIO().MousePos;
        ImVec2 local(mouse.x - p0.x, mouse.y - p0.y);   // tuval-yerel

        ImU32 col = IM_COL32((int)(draw_col[0]*255), (int)(draw_col[1]*255),
                             (int)(draw_col[2]*255), 255);

        if (ImGui::IsItemActive() && ImGui::IsMouseDown(ImGuiMouseButton_Left)) {
            if (!drawing) {
                strokes.push_back({{}, col, thickness});   // yeni stroke basla
                drawing = true;
            }
            // Ayni noktayi tekrar eklememek icin kucuk esik
            auto& pts = strokes.back().points;
            if (pts.empty() ||
                fabsf(pts.back().x-local.x)+fabsf(pts.back().y-local.y) > 2.0f)
                pts.push_back(local);
        } else {
            drawing = false;
        }

        // Cizimi kirp ve ciz (tuval disina tasmasin)
        dl->PushClipRect(p0, p1, true);
        for (auto& s : strokes) {
            for (size_t i = 1; i < s.points.size(); i++) {
                ImVec2 a(p0.x + s.points[i-1].x, p0.y + s.points[i-1].y);
                ImVec2 b(p0.x + s.points[i].x,   p0.y + s.points[i].y);
                dl->AddLine(a, b, s.color, s.thickness);
            }
        }
        dl->PopClipRect();

        ImGui::End();
    });
}
\end{lstlisting}

\notk[Yerel vs ekran koordinatı]{Noktaları \emph{tuval-yerel} olarak sakladık
(tuvalin sol-üstüne göre), çizerken \code{p0}'ı ekleyerek ekran uzayına
çevirdik. Böylece pencere taşınırsa çizim onunla birlikte gelir --- ekran
koordinatı saklasaydık pencere hareket edince çizim yerinde kalırdı. Bu
ayrım, tüm tuval/harita/node editörlerinde temel bir tasarım kararıdır.}

% =====================================================================
\gsec{Proje 7: Oyun Hata Ayıklama HUD'u}
% =====================================================================

Son projemiz, ImGui'nin en yaygın kullanım alanını temsil eder: bir oyun/render
döngüsünün üzerine bindirilen hata ayıklama arayüzü. Şeffaf bir üst-katman
(overlay), gerçek zamanlı ayar panelleri ve bir varlık (entity) denetleyicisi
içerir. Burada ``oyun'' sahte bir simülasyondur; önemli olan HUD kalıplarıdır.

\gsub{Tam Kaynak}

\begin{lstlisting}
#include "app.h"
#include <vector>
#include <cmath>

struct Entity {
    char name[32];
    float x, y;
    float vx, vy;
    bool alive = true;
};

int main() {
    App app("Oyun Debug HUD");

    // Sahte oyun durumu
    std::vector<Entity> entities;
    for (int i = 0; i < 6; i++) {
        Entity e; snprintf(e.name, sizeof(e.name), "Dusman_%d", i);
        e.x = 100 + i*90; e.y = 200;
        e.vx = cosf(i)*60; e.vy = sinf(i)*60;
        entities.push_back(e);
    }
    int selected = -1;
    bool paused = false;
    float sim_speed = 1.0f;
    bool show_overlay = true, show_inspector = true;

    app.Run([&]() {
        ImGuiIO& io = ImGui::GetIO();
        float dt = paused ? 0.0f : io.DeltaTime * sim_speed;

        // --- "Oyun" mantigi ---
        for (auto& e : entities) {
            if (!e.alive) continue;
            e.x += e.vx * dt; e.y += e.vy * dt;
            if (e.x < 0 || e.x > 1280) e.vx = -e.vx;
            if (e.y < 0 || e.y > 720)  e.vy = -e.vy;
        }

        // --- Varliklari dogrudan ekrana ciz (arka-plan DrawList) ---
        ImDrawList* bg = ImGui::GetBackgroundDrawList();
        for (int i = 0; i < (int)entities.size(); i++) {
            Entity& e = entities[i];
            if (!e.alive) continue;
            ImU32 col = (i == selected) ? IM_COL32(255,220,0,255)
                                        : IM_COL32(220,80,80,255);
            bg->AddCircleFilled(ImVec2(e.x, e.y), 14, col);
            bg->AddText(ImVec2(e.x-14, e.y+16), IM_COL32(255,255,255,200), e.name);
        }

        // --- 1) Ust menu ---
        if (ImGui::BeginMainMenuBar()) {
            if (ImGui::BeginMenu("Debug")) {
                ImGui::MenuItem("FPS Overlay", nullptr, &show_overlay);
                ImGui::MenuItem("Inspector",   nullptr, &show_inspector);
                ImGui::EndMenu();
            }
            ImGui::EndMainMenuBar();
        }

        // --- 2) Seffaf FPS overlay (sag ust kose) ---
        if (show_overlay) {
            ImGuiViewport* vp = ImGui::GetMainViewport();
            ImGui::SetNextWindowPos(
                ImVec2(vp->WorkPos.x + vp->WorkSize.x - 10, vp->WorkPos.y + 10),
                ImGuiCond_Always, ImVec2(1, 0));   // sag-ust hizala
            ImGui::SetNextWindowBgAlpha(0.35f);
            ImGuiWindowFlags f = ImGuiWindowFlags_NoDecoration
                | ImGuiWindowFlags_NoMove | ImGuiWindowFlags_NoSavedSettings
                | ImGuiWindowFlags_NoNav | ImGuiWindowFlags_AlwaysAutoResize;
            ImGui::Begin("overlay", nullptr, f);
            ImGui::Text("FPS: %.0f  (%.2f ms)", io.Framerate, 1000.0f/io.Framerate);
            int alive = 0; for (auto& e : entities) if (e.alive) alive++;
            ImGui::Text("Canli varlik: %d / %d", alive, (int)entities.size());
            ImGui::Text("Durum: %s", paused ? "DURAKLADI" : "CALISIYOR");
            ImGui::End();
        }

        // --- 3) Kontrol paneli ---
        ImGui::Begin("Simulasyon");
        if (ImGui::Button(paused ? "Devam" : "Duraklat")) paused = !paused;
        ImGui::SameLine();
        if (ImGui::Button("Sifirla")) for (auto& e : entities) e.alive = true;
        ImGui::SliderFloat("Sim Hizi", &sim_speed, 0.0f, 3.0f);
        ImGui::End();

        // --- 4) Varlik denetleyicisi (inspector) ---
        if (show_inspector) {
            ImGui::Begin("Inspector");
            if (ImGui::BeginTable("ent", 4, ImGuiTableFlags_Borders
                                  | ImGuiTableFlags_RowBg)) {
                ImGui::TableSetupColumn("Ad");
                ImGui::TableSetupColumn("Konum");
                ImGui::TableSetupColumn("Hiz");
                ImGui::TableSetupColumn("Islem");
                ImGui::TableHeadersRow();
                for (int i = 0; i < (int)entities.size(); i++) {
                    Entity& e = entities[i];
                    ImGui::TableNextRow();
                    ImGui::PushID(i);
                    ImGui::TableNextColumn();
                    if (ImGui::Selectable(e.name, selected == i,
                            ImGuiSelectableFlags_SpanAllColumns))
                        selected = i;
                    ImGui::TableNextColumn(); ImGui::Text("%.0f, %.0f", e.x, e.y);
                    ImGui::TableNextColumn(); ImGui::Text("%.0f, %.0f", e.vx, e.vy);
                    ImGui::TableNextColumn();
                    if (ImGui::SmallButton(e.alive ? "Oldur" : "Diri"))
                        e.alive = !e.alive;
                    ImGui::PopID();
                }
                ImGui::EndTable();
            }
            // Secili varligin detayli duzenlemesi
            if (selected >= 0 && selected < (int)entities.size()) {
                ImGui::SeparatorText("Secili Varlik");
                Entity& e = entities[selected];
                ImGui::InputText("Ad", e.name, sizeof(e.name));
                ImGui::DragFloat2("Konum", &e.x, 1.0f);
                ImGui::DragFloat2("Hiz", &e.vx, 1.0f);
            }
            ImGui::End();
        }
    });
}
\end{lstlisting}

\bilgi{Bu proje, ImGui'nin oyun geliştirmedeki ``süper gücünü'' özetler:
\code{GetBackgroundDrawList()} ile oyun dünyasının \emph{üzerine} çizip
(varlık etiketleri, çarpışma kutuları, yol bulma ızgaraları), aynı anda
yüzen panellerle canlı ayar yapabilirsiniz. \code{selected} indeksiyle
tablodaki seçim ile dünyadaki vurgulama otomatik senkron kalır --- çünkü ikisi
de \emph{aynı} veriyi her kare yeniden okur. Retained mode'da bu senkronu elle
sağlamak zorunda kalırdınız.}


% #####################################################################
\gpart{5}{Başvuru ve Kapanış}
% #####################################################################

% =====================================================================
\gsec{Hızlı Başvuru (Cheat Sheet)}
% =====================================================================

\gsub{Sık Kullanılan Çağrılar}

\begin{center}
\begin{tabularx}{\textwidth}{@{}>{\ttfamily\footnotesize\color{koyuGri}}l >{\raggedright\arraybackslash}X@{}}
\toprule
\rowcolor{acikMavi}\normalfont\color{anaMavi}\textbf{Çağrı} & \color{anaMavi}\textbf{İş}\\
\midrule
Begin/End & Pencere aç/kapat (End daima çağrılır)\\
\rowcolor{acikGri} Text / TextColored & Metin göster\\
Button & Düğme; tıklandıysa \code{true}\\
\rowcolor{acikGri} Checkbox / RadioButton & İki/çok durumlu seçim\\
SliderFloat / DragFloat & Sayısal değer ayarla\\
\rowcolor{acikGri} InputText & Metin girişi\\
Combo / ListBox & Açılır/liste seçimi\\
\rowcolor{acikGri} ColorEdit4 & Renk seç\\
SameLine & Sonraki widget yan yana\\
\rowcolor{acikGri} Separator / Spacing & Ayraç / boşluk\\
PushID / PopID & Benzersiz ID kapsamı\\
\rowcolor{acikGri} BeginTable / EndTable & Tablo\\
BeginChild / EndChild & Kaydırılabilir alt bölge\\
\rowcolor{acikGri} IsItemHovered / Clicked & Son öğe durumu\\
GetWindowDrawList & Özel çizim (DrawList)\\
\bottomrule
\end{tabularx}
\end{center}

\gsub{Eşleşmesi Zorunlu Çiftler}

\begin{center}
\begin{tabularx}{\textwidth}{@{}>{\raggedright\arraybackslash}X >{\raggedright\arraybackslash}X@{}}
\toprule
\rowcolor{acikMavi}\color{anaMavi}\textbf{Kural} & \color{anaMavi}\textbf{Açıklama}\\
\midrule
\code{Begin}$\leftrightarrow$\code{End} & Pencere: \code{End} her zaman (Begin false olsa da).\\
\rowcolor{acikGri} \code{BeginChild}$\leftrightarrow$\code{EndChild} & Child için de End her zaman.\\
\code{TreeNode}$\rightarrow$\code{TreePop} & Yalnızca TreeNode \code{true} dönerse.\\
\rowcolor{acikGri} \code{BeginTable}$\rightarrow$\code{EndTable} & Yalnızca Begin \code{true} dönerse.\\
\code{Push*}$\leftrightarrow$\code{Pop*} & Push sayısı = Pop sayısı (ID, StyleColor, StyleVar).\\
\bottomrule
\end{tabularx}
\end{center}

% =====================================================================
\gsec{Sık Yapılan Hatalar ve Çözümleri}
% =====================================================================

\begin{center}
\begin{tabularx}{\textwidth}{@{}>{\raggedright\arraybackslash}p{0.42\textwidth} >{\raggedright\arraybackslash}X@{}}
\toprule
\rowcolor{kirmizi!12}\color{kirmizi}\textbf{Belirti} & \color{kirmizi}\textbf{Neden / Çözüm}\\
\midrule
Döngüdeki düğmeler yanlış çalışıyor & ID çakışması. \code{PushID(i)}/\code{PopID()} kullan.\\
\rowcolor{acikGri} Değer bir kare sonra sıçrıyor & Değeri \code{static} yaptın ama iki yerde çizdin; tek kaynağa bağla.\\
Assert: ``mismatched Begin/End'' & \code{End}'i atladın ya da yanlış eşledin. Çiftleri kontrol et.\\
\rowcolor{acikGri} Türkçe harfler kutu/boş çıkıyor & Font glyph aralığına Latin Extended-A ekle, UTF-8 kaydet.\\
Pencere taşınamıyor & Konumu \code{ImGuiCond\_Always} ile zorluyorsun; \code{\_FirstUseEver} kullan.\\
\rowcolor{acikGri} Metin yazarken oyun hareket ediyor & \code{io.WantCaptureKeyboard} kontrol edilmedi.\\
Uzun liste FPS düşürüyor & \code{ImGuiListClipper} kullan.\\
\rowcolor{acikGri} Stil tüm UI'da bozuldu & \code{PushStyleColor/Var} sayısı Pop ile eşleşmiyor.\\
Ekranda hiçbir şey yok & \code{NewFrame}/\code{Render} sırası veya \code{RenderDrawData} eksik.\\
\bottomrule
\end{tabularx}
\end{center}

% =====================================================================
\gsec{Sonraki Adımlar}
% =====================================================================

Bu rehber Dear ImGui'nin çekirdeğini ve en yaygın kullanım kalıplarını
kapsadı. Buradan sonra derinleşmek için:

\begin{center}
\begin{tabularx}{\textwidth}{@{}>{\bfseries\color{anaMavi}}l >{\raggedright\arraybackslash}X@{}}
\toprule
\rowcolor{acikMavi}\color{anaMavi}\textbf{Konu / Araç} & \color{anaMavi}\textbf{Ne için}\\
\midrule
\code{imgui\_demo.cpp} & En iyi öğretmen: her widget'ın canlı örneği ve kaynağı.\\
\rowcolor{acikGri} ImPlot & Ciddi bilimsel/gerçek zamanlı grafikler.\\
ImGuizmo & 3B sahnede taşı/döndür/ölçekle gizmoları.\\
\rowcolor{acikGri} imgui-node-editor & Düğüm tabanlı editörler (shader graph vb.).\\
Docking branch & IDE benzeri çoklu-pencere düzenleri.\\
\rowcolor{acikGri} Custom backend & Kendi motorunuza (DX12/Vulkan/Metal) entegre etme.\\
\code{imgui\_stdlib.h} & \code{std::string} ile InputText.\\
\bottomrule
\end{tabularx}
\end{center}

\bilgi{ImGui öğrenmenin en hızlı yolu \textbf{yazmaktır}. Bir yan projenizi
açın, üstüne bir \code{ImGui::Begin("Debug")} penceresi koyun ve o projeye
özgü değerleri (pozisyonlar, sayaçlar, bayraklar) canlı gösterip ayarlamaya
başlayın. Bir kez bu alışkanlığı edindiğinizde, ImGui'siz geliştirme yapmak
zor gelecek. İyi çalışmalar!}

\vfill
\begin{center}
{\color{ortaGri}\rule{0.5\linewidth}{0.8pt}}\\[6pt]
{\color{ortaGri}\small C++ ile Dear ImGui Rehberi --- Son}\\[2pt]
{\color{ortaGri}\footnotesize Dear ImGui 1.90+ \textbullet\ GLFW 3.x \textbullet\ OpenGL 3.3 \textbullet\ C++17}
\end{center}

\end{document}
